\section{A. The conormal step}
\subsection{Reduced hypersurfaces and the exact kernel}

Throughout this section $i:S\hookrightarrow Y$ is the inclusion of a reduced effective Cartier divisor into a complex manifold of dimension $n+1$, $j:S_{\rm reg}\hookrightarrow S$ is the inclusion of its regular locus, and $A$ is an invertible analytic $\OO_S$-module. We write
\[
 N^*=\OO_Y(-S)|_S,\qquad V=\Omega_Y^1|_S,\qquad
 T=\ker(\Omega_S^1\longrightarrow j_*\Omega^1_{S_{\rm reg}}),\qquad Q=\Omega_S^1/T.
\]
All maps between sheaves below are $\OO_S$-linear maps of coherent analytic sheaves, except for the displayed restriction into $j_*\Omega^1_{S_{\rm reg}}$. The latter is also $\OO_S$-linear; its kernel and image have the coherent descriptions proved below.

\begin{statement}{A.1.1}{Proved}{The conormal filtration and its exact section criterion}
The conormal map $d:N^*\to V$, defined by $[F]\mapsto dF|_S$ for a local defining equation $F$, is injective. With $K=\ker(V\to Q)$, there are canonical exact sequences
\begin{align}
0&\longrightarrow N^*\xrightarrow{d}K\longrightarrow T\longrightarrow0,
\label{A:filtration1}\\
0&\longrightarrow K\longrightarrow V\longrightarrow Q\longrightarrow0.
\label{A:filtration2}
\end{align}
Let $\partial_A:H^0(S,Q\otimes A)\to H^1(S,K\otimes A)$ and $\delta_A:H^0(S,T\otimes A)\to H^1(S,N^*\otimes A)$ be their connecting maps. Then
\[
 H^0(S,V\otimes A)=0
 \quad\Longleftrightarrow\quad
 H^0(S,K\otimes A)=0\ \hbox{and}\ \partial_A\hbox{ is injective}.
\tag{A.1}
\]
Under the two vanishings in Lead 1, $H^0(S,V\otimes A)$ is canonically $\ker\delta_A$.
\end{statement}
\begin{proof}
Locally $R=\C\{x_0,\ldots,x_n\}/(F)$ is reduced. Its total quotient ring is the product of the fraction fields of its irreducible components. At the generic point of every component, $S$ is smooth and $dF$ is nonzero. Consequently $a\,dF=0$ in $R^{n+1}$ implies that $a$ vanishes in each of those fields. The map from the reduced ring $R$ into their product is injective, so $a=0$. This proves injectivity in the usual right-exact conormal sequence. Taking the inverse image of $T\subset\Omega_S^1$ gives \eqref{A:filtration1} and \eqref{A:filtration2}. Tensoring by $A$ is exact because $A$ is locally free. The cohomology sequence of \eqref{A:filtration2} identifies the image of $H^0(VA)$ in $H^0(QA)$ with $\ker\partial_A$ and its kernel with $H^0(KA)$, proving (A.1). When $H^0(QA)=0$, it identifies $H^0(VA)$ with $H^0(KA)$. The first sequence then identifies this group with $\ker\delta_A$ when $H^0(N^*A)=0$.
\end{proof}

\begin{statement}{A.2.1}{Proved}{The Jacobian fractional ideal for every reduced hypersurface}
Let $\mathcal J\subset\OO_S$ be the Jacobian ideal, generated locally by $F_{x_0},\ldots,F_{x_n}$ modulo $F$. In the total quotient sheaf $\mathcal M_S$ define
\[
 \mathcal J^{-1}=\{q\in\mathcal M_S:q\mathcal J\subset\OO_S\}.
\]
There are canonical identifications
\[
 K\simeq N^*\otimes\mathcal J^{-1},\qquad
 T\simeq N^*\otimes(\mathcal J^{-1}/\OO_S),
\tag{A.2}
\]
under which the first map in \eqref{A:filtration1} is induced by $\OO_S\hookrightarrow\mathcal J^{-1}$. Thus (A.1), with this explicitly defined fractional ideal, is an exact necessary and sufficient criterion for every reduced hypersurface.
\end{statement}
\begin{proof}
The ideal $J$ contains a nonzerodivisor: it is nonzero at every minimal prime, and prime avoidance supplies an element outside all minimal primes. Every $R$-linear map $\varphi:J\to R$ is multiplication by a unique element of the total quotient ring. Indeed, for a nonzerodivisor $a\in J$, the identity $a\varphi(b)=b\varphi(a)$ for $b\in J$ gives the multiplier $q=\varphi(a)/a$. Hence $J^{-1}=\Hom_R(J,R)$ is coherent.

At the generic points, an element $\alpha\in K$ is uniquely $q\,dF$ for a tuple of meromorphic functions $q$. The coefficients of $\alpha$ belong to $R$ exactly when $qF_{x_i}\in R$ for every $i$, which is the condition $q\in J^{-1}$. Conversely such a vector restricts to zero in the cotangent space of every regular component, hence belongs to $K$. The map $[F]\otimes q\mapsto q\,dF|_S$ proves the first identification. Its quotient by $R\,dF$ gives the second. Replacing $F$ by $uF$, with $u$ a holomorphic unit, replaces $dF|_S$ by $u\,dF|_S$ and replaces the conormal generator by $u[F]$. Coordinate changes transform the derivative vector by an invertible Jacobian matrix. These are exactly the transitions in (A.2), so the identifications are canonical.
\end{proof}

\begin{statement}{A.2.2}{Proved}{Normalization, conductor, and the excess kernel}
Let $\eta:\widetilde S\to S$ be the finite analytic normalization, $\mathcal B=\eta_*\OO_{\widetilde S}$, and $\mathfrak c=\Ann_{\OO_S}(\mathcal B/\OO_S)$ its conductor. There are canonical inclusions
\[
 \mathcal J\subset\mathfrak c,\qquad
 \mathcal B\subset\mathcal J^{-1},\qquad
 \eta_*\eta^*N^*\hookrightarrow K.
\tag{A.3}
\]
In a local ring $R$ with normalization $B$, put $D=\Hom_B(JB,B)$, viewed as a fractional ideal. Then
\[
 J^{-1}=\ker\left[D\longrightarrow
 \Hom_R\bigl(J,(B/\mathfrak c)/(R/\mathfrak c)\bigr)\right],
\quad q\longmapsto\bigl(b\mapsto qb\bmod R\bigr).
\tag{A.4}
\]
For any locally free $\OO_S$-module $E$, the conductor gives a second, independent exact section criterion:
\begin{multline}
H^0(S,E)=
\ker\left[H^0(\widetilde S,\eta^*E)\oplus H^0(D_c,E|_{D_c})\right.\\
\left.\longrightarrow H^0(\widetilde D_c,\eta^*E|_{\widetilde D_c})\right],
\qquad (a,b)\longmapsto a|_{\widetilde D_c}-\eta^*b,
\tag{A.5}
\end{multline}
where $D_c$ and $\widetilde D_c$ are defined by $\mathfrak c$ in $S$ and $\widetilde S$. In particular, the ambient group vanishes exactly when the map in (A.5), for $E=V\otimes A$, is injective.
\end{statement}
\begin{proof}
We prove $JB\subset R$ rather than assume a conductor formula. Choose a generic linear projection for which Weierstrass preparation writes the equation as a unit times a monic polynomial $P(z)$ over $A_0=\C\{x_1,\ldots,x_n\}$. Thus $R=A_0[z]/(P)$ is finite free over $A_0$ and generically separable. The functional taking the coefficient of $z^{m-1}$ in a remainder modulo $P$, where $m=\deg P$, generates $\Hom_{A_0}(R,A_0)$ as an $R$-module: the matrix of the pairing on $1,z,\ldots,z^{m-1}$ has zeros below the reverse diagonal and ones on that diagonal. Over the fraction field, Lagrange interpolation identifies that functional with $r\mapsto\operatorname{Tr}(r/P'(z))$. Therefore the trace dual is $P'(z)^{-1}R$.

For $b\in B$ and $r\in R$, the element $br$ is integral over the integrally closed ring $A_0$. Its trace belongs to $A_0$. It follows that $b$ lies in this trace dual, so $P'(z)b\in R$. Modulo $F$, the derivative in the chosen projection direction differs from $P'(z)$ by a unit. The permissible projection directions form a nonempty open set; choosing a linear basis of such directions proves $F_{x_i}B\subset R$ for every $i$. This proves $J\subset\mathfrak c$ and hence $B\subset J^{-1}$. The last inclusion of (A.3) is (A.2) and the projection formula.

The module $D$ consists exactly of the meromorphic $q$ satisfying $qJ\subset B$. Requiring that these products descend to $R$ is precisely the vanishing of their classes in $B/R=(B/\mathfrak c)/(R/\mathfrak c)$, proving (A.4). Finally the ring map
\[
 R\longrightarrow B\oplus R/\mathfrak c,
 \qquad r\longmapsto(r,r\bmod\mathfrak c)
\]
has image the kernel of $(b,\bar r)\mapsto b-\bar r\in B/\mathfrak c$. Surjectivity onto $B/\mathfrak c$ comes from the first summand. Tensoring this exact sequence by the free module corresponding to $E$ and taking global sections gives (A.5). All conductor restrictions are scheme-theoretic, so this argument includes nonreduced conductors.
\end{proof}

The fractional ideal appearing in (A.2) has an established interpretation. Granger--Schulze, Proposition~3.4 and Corollary~3.5 of \cite{GS2014}, identify $J^{-1}$ with the module of logarithmic residues and prove the inclusions in (A.3). The preceding proof gives the maps needed here directly. Their Theorems~1.1 and~1.3 also give the following exact extension of the normal-crossings case.

\begin{statement}{A.2.3}{Known, with derivation}{The precise range of the normalization formula}
The canonical inclusion $\eta_*\eta^*N^*\hookrightarrow K$ is an isomorphism exactly when $S$ has normal crossings outside an analytic subset of codimension at least two in $S$.
\end{statement}
\begin{proof}
Locally a logarithmic one-form can be written as $\alpha/F$, with $\alpha$ holomorphic and $\alpha\wedge dF$ divisible by $F$. Restriction to $S$ is uniquely $q\,dF$ at its generic points; the holomorphic coefficients of $\alpha$ give $q\in J^{-1}$. Conversely, for $q\in J^{-1}$ choose holomorphic lifts of the coefficients $qF_{x_i}$. Their exterior product with $dF$ vanishes modulo $F$, so they define a logarithmic form with residue $q$. The kernel consists exactly of holomorphic forms. Thus the logarithmic residue module in \cite[Proposition~3.4]{GS2014} is exactly our $J^{-1}$, with no change of fractional-ideal embedding. Weakly holomorphic functions are exactly $B$. Theorems~1.1 and~1.3 there say that all logarithmic residues lie in $B$ exactly when the stated normal-crossings condition holds. Combining this with $B\subset J^{-1}$ and (A.2) proves the result.
\end{proof}

\subsection{Double curves, triple points, pinch points, and cusps}

\begin{statement}{A.1.2}{Proved}{The canonical normal-crossings formula}
For a reduced normal-crossings divisor, the map in (A.3) is an isomorphism. In a chart with $F=x_1\cdots x_r$, it is explicitly
\[
 \bigoplus_{i=1}^r R/(x_i)\longrightarrow K,
 \qquad (a_i)_i\longmapsto
 \sum_{i=1}^r a_i\prod_{j\ne i}x_j\,dx_i.
\tag{A.6}
\]
The tensor factor $N^*$ is trivialized here by $[F]$. At a double crossing $T\simeq R/(x,y)$; at a triple crossing,
\[
 T\simeq R^3/\bigl(R(1,1,1)+R(x,0,0)+R(0,y,0)+R(0,0,z)\bigr).
\tag{A.7}
\]
\end{statement}
\begin{proof}
On the normalization branch $x_i=0$, an ambient form maps to zero in the cotangent bundle of that branch precisely when all coefficients except that of $dx_i$ vanish. Thus a coefficient of $dx_i$ must vanish on every branch $x_j=0$ for $j\ne i$. In the convergent power-series ring it is divisible by $\prod_{j\ne i}x_j$. Its remaining factor is defined modulo $x_i$, which gives exactly (A.6). The coefficients of different $dx_i$ show injectivity. Coefficients of the coordinates not among $x_1,\ldots,x_r$ vanish on every branch and hence vanish in $R$. This proves surjectivity in every dimension.

For $r=1$ the map is $a\mapsto a\,dx$. For $r=2$ it is $(a,b)\mapsto ay\,dx+bx\,dy$, and quotienting by $dF=y\,dx+x\,dy$ leaves the generator $y\,dx=-x\,dy$, annihilated by $(x,y)$. For $r=3$ the branch generators are $yz\,dx,xz\,dy,xy\,dz$. Their sum is $dF$ and their respective annihilators are $(x),(y),(z)$, proving (A.7). Multiplication of $F$ by a unit and changes of coordinates were checked in A.2.1. Formula (A.6) is the restriction of the canonical fractional-ideal map, so it glues without a chosen branch ordering, equation, or coordinate system.
\end{proof}

\begin{statement}{A.2.4}{Proved}{Pinch points and cuspidal edges}
At the pinch point $F=x^2-y^2z$, with normalization $x=uv,y=u,z=v^2$, write $B=\C\{u,v\}$. Then
\[
 K=B\,dF,\qquad T\simeq B/R\simeq R/(x,y),
\quad \tau=2yz\,dx-2xz\,dy-xy\,dz
\tag{A.8}
\]
is a torsion generator. At the cuspidal edge $F=x^2-y^3$, with normalization $x=t^3,y=t^2$ and unchanged coordinate $z$, write $B=\C\{t,z\}$. Then
\[
 K=t^{-1}B\,dF,\qquad
 T\simeq t^{-1}B/R\simeq R/(x,y^2),
 \quad\tau=2y\,dx-3x\,dy.
\tag{A.9}
\]
The conductor is $(x,y)$ for the pinch point and $(x,y)$ for the cusp; in the cusp normalization it is $t^2B$. An ordinary surface, meaning a surface whose singularities are double normal crossings, triple normal crossings, and pinch points, satisfies canonically $K=\eta_*\eta^*N^*$ and $T=N^*\otimes(\eta_*\OO_{\widetilde S}/\OO_S)$. For the cuspidal edge the additional quotient $J^{-1}/B$ is generated by $t^{-1}$ over $\C\{z\}$.
\end{statement}
\begin{proof}
For the pinch point,
\[
 R=\C\{u,v^2,uv\}=\C\{v^2\}+uB,
 \quad J=(uv,uv^2,u^2)=u(u,v)B.
\]
The equality of the two descriptions of $J$ follows by separating even and odd powers of $v$. An element $q\in J^{-1}$ satisfies $(qu)(u,v)\subset B$. Since $B$ is a regular local ring, $\Hom_B((u,v),B)=B$: a meromorphic function whose products with $u$ and $v$ are holomorphic has no pole along any divisor, and its possible isolated singularity is removable by the holomorphic extension theorem. Hence $qu=h\in B$. The requirements $hv,hv^2\in R$, restricted to $u=0$, say respectively that $h(0,v)v$ and $h(0,v)v^2$ are even power series. The first requires $h(0,v)$ to be odd and the second to be even. Therefore $h(0,v)=0$, so $h\in uB$ and $q\in B$. The reverse inclusion follows either directly or from (A.3). Also $B/R$ is generated by $v$ and its annihilator is $(u,uv)=(y,x)$. Multiplying $dF$ by $v=x/y$ gives the displayed $\tau$.

For the cusp, $R=\C\{t^2,t^3,z\}$ and $J=(t^3,t^4)$. A meromorphic $q$ with $qt^3,qt^4\in R$ has no poles away from $t=0$; writing its Laurent expansion in $t$ shows that the lowest allowed exponent is $-1$. Indeed an exponent $-3$ produces a forbidden $t$ term after multiplication by $t^4$, an exponent $-2$ produces one after multiplication by $t^3$, and smaller exponents produce a pole. Conversely every $t^{-1}B$ has both products in $R$. Thus $J^{-1}=t^{-1}B$. Modulo $R$, the quotient has the $\C\{z\}$-basis $t^{-1},t$; multiplication by $y$ sends the first basis vector to the second, while $x$ and $y^2$ kill both. This proves the cyclic presentation and (A.9). The conductor calculations follow directly from the displayed descriptions of $R\subset B$. Formula (A.2) and the three local computations prove the ordinary-surface assertion by gluing.
\end{proof}

\begin{statement}{A.2.5}{Proved}{Vertical residue poles at a tacnode--cusp intersection}
For $F=x^2-y^4z^3$, normalization is $x=u^2v^3$, $y=u$, $z=v^2$. Put $A_0=\C\{u,v^2\}$ and $B=\C\{u,v\}$. Then
\[
 \begin{aligned}
 R&=A_0\oplus u^2v^3A_0,&\mathfrak c&=u^2v^2B,\\
 J^{-1}&=A_0\oplus(uv)^{-1}A_0,&
 J^{-1}/B&\simeq A_0/(uv^2).
 \end{aligned}
 \tag{A.9a}
\]
In a trivialization of $N^*$ the torsion module is
\[
 T\simeq R/(x,y^3z^2),\qquad
 \tau=2yz\,dx-4xz\,dy-3xy\,dz.
 \tag{A.9b}
\]
In particular $v^{-1}\in J^{-1}\setminus B$. For the ruling with base coordinate $v$, this is a residue constant on its generic fibres with a pole on the vertical divisor $v=0$, although the horizontal conductor $u=0$ has coefficient two.
\end{statement}
\begin{proof}
The element $v=x/(y^2z)$ belongs to the fraction field and satisfies $v^2=z$. Adjoining it gives the regular ring $B$, which is therefore the finite normalization. Separating even and odd powers of $v$ gives the displayed description of $R$. An even monomial $u^av^b$ belongs to the conductor exactly when $a\ge2$ and $b\ge2$; an odd monomial already in $R$ has $a\ge2$, $b\ge3$ and is then in the conductor. Thus $\mathfrak c=u^2v^2B$.

The normalized Jacobian generators, with irrelevant nonzero constants removed, are $u^2v^3,u^3v^6,u^4v^4$. If $qJ\subset R$, the first product shows that $q$ has no other polar divisor, with possible orders at most two along $u=0$ and three along $v=0$. A term of order $v^{-3}$ would produce a forbidden $v^1$ term in $qu^4v^4$; a term of order $v^{-2}$ would produce one in $qu^2v^3$. Both coefficients must vanish, so $q$ has order at least $-1$ in $v$. A coefficient $a(v)u^{-2}$ would require both $a(v)v^3$ and $a(v)v^6$ to be even Laurent series, from the $u^0$ and $u^1$ coefficients of the first two products. The first requirement makes $a$ odd and the second makes it even, so $a=0$. Thus $q=t/(uv)$ with $t\in B$.

Write $t=a(u,v^2)+v b(u,v^2)$. The first Jacobian product equals $uv^2t$ and belongs to $R$ exactly when $b$ is divisible by $u$. The other two products automatically belong to $R$ under this condition. This gives $J^{-1}=A_0\oplus(uv)^{-1}A_0$. The odd part of $B$ is $(uv)^{-1}(uv^2)A_0$, giving $J^{-1}/B$ in (A.9a). The odd part of $R$ is $(uv)^{-1}u^3v^4A_0$, so $J^{-1}/R$ is generated by $(uv)^{-1}$ with annihilator $(x,y^3z^2)$. Multiplication of $dF$ by $(uv)^{-1}$ gives (A.9b).

The particular vertical pole has the completely regular gradient products
\[
 v^{-1}dF=2y^2z\,dx-4xyz\,dy-3xy^2\,dz.
\]
Pulling this expression to $B$ gives zero. It is a nonzero ambient-kernel vector, with multiplier $v^{-1}$, rather than a section of the normalization submodule. The exact identities and finite monomial checks are saved in \path{code/A_vertical_pole.py}; the Laurent argument above establishes the full convergent module. This local example does not assert the existence of a compact CDP fibration realizing that configuration.
\end{proof}

For ordinary surfaces the exact ambient-vanishing criterion is (A.1) with $K\otimes A=\eta_*\eta^*(N^*\otimes A)$. For surfaces with cuspidal-edge charts it is (A.1) with the locally specified sheaf $N^*\otimes t^{-1}B$. These specify the first group and the target of the connecting homomorphism, including the extra polar term at the cusp. Under $H^0(QA)=0$ they reduce respectively to the vanishing of these explicitly computed kernel sections. The general criterion (A.4)--(A.5) retains the full conductor and applies without a singularity-type restriction. The exact polynomial identities and exponent computations are recorded in \path{code/A_local_algebra.py} and its output.

\begin{statement}{A.2.6}{Proved}{A transverse cusp whose critical-curve equation is only a symbolic square}
Put
\[
 \begin{aligned}
 I&=(y^2-xz,\ yz^2-x^4,\ z^3-x^3y)\subset\C\{x,y,z\},\\
 F&=4x^7-15x^3yz^2+10x^2y^3z-3xy^5+4z^5.
 \end{aligned}                                                   \tag{A.NL.1}
\]
The reduced critical locus of the reduced hypersurface $F=0$ is the
irreducible curve $D$ parametrized by $(x,y,z)=(t^5,t^6,t^7)$.
Its transverse type at every point with $t\ne0$ is the ordinary cusp
$A_2$. The curve is not a local complete intersection at the origin,
and
\[
             F\in I^{(2)}\setminus I^2.                         \tag{A.NL.2}
\]
Thus the ordinary-square membership used to define the conormal
quadratic jet in A.5.5 does not extend to an arbitrary reduced
critical curve, even with generic transverse type $A_2$.
This is a local example, without an asserted compact CDP realization.
\end{statement}
\begin{proof}
Write the three displayed generators of $I$ as $g_1,g_2,g_3$.
The quotient is a finite $\C\{x\}$-module generated by
$1,y,z,yz,z^2$: their products by $y$ and $z$ reduce to this list
using $y^2=xz$, $yz^2=x^4$, and $z^3=x^3y$.
For example these identities also give $y^5=x^6$ and $z^5=x^7$,
which make the analytic quotient finite over $\C\{x\}$.
The parametrization sends the five generators to
$1,t^6,t^7,t^{13},t^{14}$. Their exponents have distinct residues
modulo five, so they are linearly independent over $\C\{t^5\}$.
Consequently this actual ring homomorphism is injective and has
image $\C\{t^5,t^6,t^7\}$; it proves the stated normalization and
the primeness of $I$. The three generators are minimal. Indeed their
initial terms have degrees two, three, three; after the degree-two
coefficient vanishes, the cubic terms $yz^2$ and $z^3$ are independent
modulo the span of $xg_1,yg_1,zg_1$. A height-two ideal with three
minimal generators is not a complete intersection.

The following polynomial identities determine the critical locus:
\[
 \begin{aligned}
 F_y&=-15xg_1^2,\\
 F_x&=(17xyz-3y^3)g_1-28x^2g_2,\\
 F_z&=10x^2yg_1+20zg_3.
 \end{aligned}                                                   \tag{A.NL.3}
\]
At $x=0$, the equations $F_x=F_z=0$ give $y=z=0$.
At $x\ne0$, $F_y=0$ gives $g_1=0$, and then $F_x=0$
gives $g_2=0$. The value $z=0$ is impossible there, so $F_z=0$
also gives $g_3=0$. Conversely all three partials vanish on $D$.
Weighted Euler's identity, with weights $(5,6,7)$ and degree $35$,
puts $F$ in its Jacobian ideal. The critical locus is therefore
exactly $D$ as a reduced set. A repeated factor of $F$ would put
a two-dimensional component in this critical locus, which proves
that $F$ is reduced.

The Hessian along the normalization has the exact factorization
\[
 \operatorname{Hess}(F)|_D
 =2t^{21}
 \begin{pmatrix}7t^2\\0\\-5\end{pmatrix}
 \begin{pmatrix}7t^2&0&-5\end{pmatrix},
 \qquad \frac{1}{6}F_{yyy}|_D=-20t^{17}.                       \tag{A.NL.4}
\]
For $t\ne0$, the plane $x=t^5$ is transverse to $D$.
The restricted Hessian has rank one and its radical is the
$y$ direction; its cubic coefficient in that direction is nonzero.
The holomorphic splitting lemma and a local cube root of a unit
give the transverse cusp equation. At the generic point of $D$,
its conormal is locally free and injects into ambient differentials.
Since $F$ and $dF$ vanish there, $F\in I^2\C\{x,y,z\}_I$,
which is precisely $F\in I^{(2)}$.

To decide ordinary-square membership, use the positive weights
$(5,6,7)$. The three generators have weights $12,20,21$.
Among products of two generators, only $g_1^2$ can contribute
to weight $35$: it needs a coefficient of weight $11$, whose
space is spanned by $xy$. Products $g_1g_2$ and $g_1g_3$ would
need weights $3$ and $2$, which do not occur; the other products
already have greater weight. Thus the weight-$35$ part of $I^2$
is exactly $\C\,xyg_1^2$. Its $x^7$ coefficient is zero, whereas
that of $F$ is four. Taking the weight-$35$ component of a
convergent series identity is legitimate because all weights are
positive. This proves $F\notin I^2$ in the analytic local ring.

For comparison of the curve invariants, the normalization conductor
is $t^{10}\C\{t\}$: the semigroup gaps are $1,2,3,4,8,9$.
The differential image in $\C\{t\}\,dt$ is $t^4\C\{t\}\,dt$.
The smallest valuation of a two-by-two minor of the displayed
generators' Jacobian matrix is $19$. Thus the normalized conormal
determinant correction is $19-4=15$, while the conductor exponent
is $10$; the complete-intersection adjunction comparison cannot
identify these two quantities here.

The symbolic-square correction is visible directly on the free
normalized conormal. Set $a=dg_1|_D$ and
$b=dg_2|_D+2t^8a$. Then
\[
 \begin{aligned}
 a&=(-t^7,2t^6,-t^5),& b&=(-6t^{15},5t^{14},0),\\
 dg_3|_D&=-3t^9a+tb,&
 q&=t^{-5}(5t^8a-2b)^2.
 \end{aligned}                                                   \tag{A.NL.5}
\]
The first two rows generate the entire normalized conormal image,
and $q$ is the quadratic term, namely half the Hessian.
If $A,B$ are the dual basis, its radical frame $2A+5t^8B$
is represented in the ambient normal quotient by $t^{-6}\partial_y$.
Its intrinsic cubic is therefore $-20t^{-1}$. These are actual
quadratic and cubic poles of orders five and one. The corresponding
combined invariant $q^3k^2$ has pole order $17$; no bound for such
poles on all non-lci curves is deduced from this one computation.
All substitutions, polynomial identities, weighted coefficient
spaces, conormal comparisons, and Jacobian minors are checked
exactly in \path{code/A_non_lci_cusp.py}.
\end{proof}

\subsection{The full nonreduced fibre}

\begin{statement}{A.3.1}{Proved}{All thickenings and their lifting maps}
Let $S$ be any effective Cartier divisor in a complex manifold, let $E$ be any locally free sheaf on that manifold, and let $E_k=E|_{kS}$. For each $k\ge2$ there is a canonical exact sequence
\[
0\longrightarrow (E|_S)\otimes(N^*)^{k-1}
 \longrightarrow E_k\longrightarrow E_{k-1}\longrightarrow0.
\tag{A.10}
\]
Write $G_{k-1}=(E|_S)\otimes(N^*)^{k-1}$ and let
$\epsilon_k:H^0((k-1)S,E_{k-1})\to H^1(S,G_{k-1})$
be the connecting map. Then
\[
0\longrightarrow H^0(S,G_{k-1})\longrightarrow H^0(kS,E_k)
 \longrightarrow\ker\epsilon_k\longrightarrow0.
\tag{A.11}
\]
In particular $H^0(kS,E_k)=0$ exactly when $H^0(S,G_{k-1})=0$ and $\epsilon_k$ is injective. Starting with $H^0(S,E|_S)$, equations (A.10)--(A.11) compute the group on $mS$ through all its extension maps. The vanishing of every $H^0(S,G_j)$, $0\le j<m$, is a sufficient criterion. These statements apply to $E=\Omega_Y^1\otimes B$ for every line bundle $B$, and include smooth, normal-crossings, and all other reduced supports.
\end{statement}
\begin{proof}
With $I=\OO_Y(-S)$, the sequence is the tensor product of $E$ with
$0\to I^{k-1}/I^k\to\OO_Y/I^k\to\OO_Y/I^{k-1}\to0$.
The Cartier property identifies $I^{k-1}/I^k$ with $(N^*)^{k-1}$ on $S$; local freeness makes the tensor product exact. Taking global sections proves (A.11) and its equivalence. Induction on $k$ proves the sufficient criterion. No splitting of the thickening, cyclic base change, or assumption about the normal bundle has been used.
\end{proof}

\begin{statement}{A.3.2}{Refuted, with computed replacement}{Vanishing on the support alone does not imply vanishing on a multiple fibre}
There is a compact smooth complex threefold $g:Y\to\PP^1$ with connected general fibre an abelian surface and a smooth fibre support $S$ of multiplicity $m$, for every $m\ge3$, and a line bundle $B$ on $Y$, such that
\[
H^0(S,(\Omega_Y^1\otimes B)|_S)=0,
\qquad h^0(mS,(\Omega_Y^1\otimes B)|_{mS})=3.
\tag{A.12}
\]
Its conormal has exact order $m$. The replacement is (A.10)--(A.11), with the graded twists retained.
\end{statement}
\begin{proof}
Take an abelian surface $F$ and a translation $a:F\to F$ of exact order $m$. Let $\zeta$ be a primitive $m$th root of unity. The action generated by $(p,z)\mapsto(a(p),\zeta z)$ on $F\times\PP^1$ is free, so the quotient $Y$ is a smooth compact threefold and the quotient map is finite etale. The invariant map $t=z^m$ gives $g:Y\to\PP^1$. Away from $0$ and $\infty$ its fibre is isomorphic to $F$: the $m$ lifts of a base point are permuted transitively and freely. At $0$ the fibre is $mS$, where $S=F/\langle a\rangle$. Set $B=\OO_Y(-S)$.

The connected cyclic cover $F\to S$ identifies $N^*=B|_S$ with the character line of weight one. Its order is exactly $m$: otherwise a nontrivial proper character power would have a nonzero invariant trivialization on the connected compact cover, whose holomorphic functions are constant. Pullback to $F$ splits the ambient cotangent bundle into two invariant tangent forms and the normal form $dz$ of weight one. Consequently
\[
 (\Omega_Y^1\otimes B)|_S\simeq (N^*)^{\oplus2}\oplus(N^*)^2.
\]
The nontrivial flat lines on the abelian surface $S$ have no holomorphic sections. For completeness, their pullbacks to $F$ are trivial; a section would be a constant function transforming by a nontrivial character, which is zero. Since $m\ge3$, both characters in the preceding splitting are nontrivial, proving the first equality.

The inverse image of $mS$ is the ringed space with structure sheaf $\OO_F\otimes\C[z]/(z^m)$. Write $e$ for the local generator of the pulled-back line $B$; it has weight one. All global sections on this inverse image are linear combinations of $z^k\alpha_1e,z^k\alpha_2e,z^kdz\,e$ for $0\le k<m$, where $\alpha_1,\alpha_2$ are a basis of invariant one-forms on $F$. The first two have weight $k+1$ and the third weight $k+2$. The invariant monomials are exactly
$z^{m-1}\alpha_1e,z^{m-1}\alpha_2e,z^{m-2}dz\,e$.
Taking invariants is exact over $\C$, by averaging over the finite group. These three monomials therefore give the claimed basis and prove the second equality. Translations act identically on $H^q(F,\OO_F)$, whose dimensions are $\binom{2}{q}$. The same calculation gives all fibre groups: $h^q(S,E|_S)=(0,0,0)$ and $h^q(mS,E|_{mS})=(3,6,3)$ for $q=0,1,2$. For completeness, on the compact threefold itself the two tangent summands pull back to $\OO_F\boxtimes\OO_{\PP^1}(-1)$ and have no cohomology. The normal summand is $\OO_F\boxtimes\Omega^1_{\PP^1}(-\{0\})$. Its first cohomology on $\PP^1$ has the Cech basis $dz/z,dz$, of weights zero and one, respectively. Indeed sections near zero are multiples of $z\,dz$, and sections near infinity are multiples of $z^{-2}dz$, leaving exactly those two Laurent classes. Its invariant part is one-dimensional. Kunneth and finite-cover descent therefore give $h^q(Y,E)=(0,1,2,1)$ for $q=0,1,2,3$.
\end{proof}

For this example the complete dimensions on $kS$, $1\le k\le m$, are $0$ for $k\le m-2$, $1$ for $k=m-1$, and $3$ for $k=m$; the same character calculation proves this. More generally the character count with $B=\OO_Y(-bS)$ is
\[
\begin{aligned}
h^0(kS,(\Omega_Y^1\otimes B)|_{kS})
&=2\#\{0\le j<k:j+b\equiv0\pmod m\}\\
&\quad+\#\{0\le j<k:j+b+1\equiv0\pmod m\}.
\end{aligned}
\tag{A.13}
\]
The finite counts for $3\le m\le12$ and all characters are checked in \path{code/A_multiple_fibres.py}.

\begin{statement}{A.3.3}{Proved}{The precise finite-translation step in Remarks 7.2 and 7.4}
For a fixed multiple fibre $mS$, suppose that the set of line bundles $B$ for which $H^0(S,(\Omega_Y^1\otimes B)|_S)=0$ contains a dense open subset $U$ of the relevant Picard component. The intersection
\[
\bigcap_{j=0}^{m-1}\{B:B\otimes\OO_Y(-jS)\in U\}
\tag{A.14}
\]
is a dense open subset on which the full-fibre ambient group vanishes, whenever the displayed translations preserve that component. The same assertion for a finite union of components follows by using the appropriate dense open subset in each translated component.
\end{statement}
\begin{proof}
Tensoring by the fixed line bundle $\OO_Y(-jS)$ is a biholomorphic translation between Picard components. Its inverse image of a dense open set is dense and open. The finite intersection retains these properties. Its defining condition is exactly the vanishing of every graded group in A.3.1, which proves the result.
\end{proof}

This finite-translation argument supplies the thickening step intended in Remark~7.2 once ambient vanishing on the smooth support is established. On a smooth support, the conormal sequence has $Q=\Omega_S^1$ and $K=N^*$, so the two requisite support vanishings do establish that ambient vanishing. Remark~7.4 requires the same argument on a singular support; a statement solely about $Q=\widetilde\Omega_S^1$ leaves the kernel in A.1.1 to be checked. Thus the thickening operation itself has a complete exact treatment, while the cited singular-support assertion cannot supply its missing kernel hypothesis. Example A.3.2 disproves the inference from a single support vanishing even with torsion conormal; it does not refute a generic-vanishing assertion already established for all the finitely translated twists.

\begin{statement}{A.3.4}{Proved}{The smooth-multiple assertion of Remark 7.2}
Under the hypotheses of CDP20 Proposition~7.3, let $mS$ be a full fibre whose reduced support $S$ is smooth. For every real irrational transgression parameter $r$ of D.1.9,
\[
 H^0\bigl(mS,(\Omega_Y^1\otimes L_r)|_{mS}\bigr)=0.
 \tag{A.15a}
\]
The same simultaneous vanishing on all such fibres holds on the complement of a proper complex analytic subset of $\Pic(Y)$. The assertion includes the entire scheme-theoretic fibre, with every nilpotent layer.
\end{statement}
\begin{proof}
Choose a disc $\Delta$ about the fibre value and write $U=g^{-1}(\Delta)$. Form the base change $t=s^m$ and normalize it, obtaining $p:U'\to U$ and $g':U'\to\Delta'$. This map $p$ is finite étale of degree $m$. Indeed, near a point of the smooth support write $g=u z^m$ with $u$ a holomorphic unit. The normalization is given by
\[
 w^m=u,\qquad s=wz.
\]
The derivative $m w^{m-1}$ is invertible, so these equations give an unramified cover of the original smooth chart. They cover the normalized base change, since they coincide away from $z=0$ and are normal. The local constructions glue by uniqueness of normalization. The divisor $S'=(g')^{-1}(0)$ is smooth and reduced, $g'$ is a submersion near it, and $p^{-1}(mS)=mS'$ as Cartier divisors.

Pull back the unitary local system (A.22) on $U$ to $U'$; its associated holomorphic line is $A'=p^*L_r$. Every nearby smooth fibre of $g'$ maps isomorphically onto the corresponding nearby fibre of $g$. Its character is therefore the irrational exponential of a nonzero integral class. The smooth-fibre calculation in A.4.9 gives zero cohomology in degrees zero and one for that local system. Proper smooth local triviality at $S'$ transports the same cohomology statement to $S'$. If a component presentation of the cover is used, the argument applies on each component and the groups add.

A holomorphic section of a unitary-flat line on a compact smooth complex manifold is parallel: its squared norm is plurisubharmonic and hence constant by the maximum principle. Thus $H^0(S',A')=0$ follows from the zero degree-zero local-system cohomology. A.4.7 gives $H^0(S',\Omega_{S'}^1\otimes A')=0$ from degree-one vanishing. The conormal of the reduced full fibre $S'$ is trivial, with global generator $dg'$. Its conormal sequence, tensored by $A'$, now gives
$H^0(S',(\Omega_{U'}^1\otimes A')|_{S'})=0$.

All the graded factors for $mS'$ in A.3.1 equal this same restricted bundle, because the conormal is trivial. Induction in (A.10) proves $H^0(mS',(\Omega_{U'}^1\otimes A')|_{mS'})=0$. Since $p$ is étale, $p^*\Omega_U^1=\Omega_{U'}^1$, and finite faithfully flat pullback injects sections on $mS$ into sections on $mS'$. One can verify that injection locally by tensoring the inclusion of a module into its faithfully flat scalar extension, or by the trace divided by $m$. This proves (A.15a).

There are finitely many non-submersion values. For each of these fixed full fibres the coherent ambient bundle varies in the holomorphic Picard family constructed in A.4.9, and proper semicontinuity makes its nonvanishing locus analytic. Each locus is proper because it excludes every $L_r$ with $r$ real irrational. Their finite union is proper, proving the simultaneous generic assertion. This proves the smooth-multiple conclusion of Remark~7.2. For a singular support the identical root cover retains its singularities and the kernel $N^*\otimes\mathcal J^{-1}$; the present proof does not erase that term in Remark~7.4.
\end{proof}

\begin{statement}{A.3.5}{Proved}{Every reducible full fibre has generic ambient vanishing}
Let $Y$ be a smooth compact complex threefold with $H^2(Y,\Q)=0$, and let $F=\sum_i m_iS_i$ be an effective Cartier divisor with connected reducible support. Every connected component of $\Pic(Y)$ contains a dense open subset, with proper closed analytic complement, on which
\[
 H^0(F,(\Omega_Y^1\otimes L)|_F)=0.
 \tag{A.15b}
\]
The multiplicities may be unequal. For a proper connected fibration over a compact curve, the vanishing is simultaneous on all reducible full fibres after intersecting finitely many such open sets. Consequently a full fibre with ambient sections for every Picard twist has irreducible support.
\end{statement}
\begin{proof}
Every component $S_i$ meets a different component $S_j$, because the support is connected and reducible. The canonical section of $\OO_Y(S_j)$ restricts to a nonzero section on $S_i$. It is a nonzerodivisor in every local ring of $S_i$: otherwise one local branch of $S_i$ would be contained in $S_j$, and analytic continuation on the global irreducible component would give $S_i=S_j$. Its zero scheme is therefore a nonzero effective Cartier divisor $D$ on $S_i$.

For any fixed locally free sheaf $E$ on the reduced irreducible compact surface $S_i$, the injections
\[
 H^0(S_i,E(-nD))\hookrightarrow H^0(S_i,E)
\]
form a decreasing sequence of finite-dimensional subspaces. Their intersection is zero. At a point of $D$, a section in the intersection has its germ in every power of the Cartier ideal; the Krull intersection theorem for the Noetherian analytic local ring makes this germ zero. The identity theorem on the reduced irreducible surface, applied in local trivializations of $E$, then makes the global section zero. A descending sequence of finite-dimensional subspaces stabilizes, so $H^0(S_i,E(-nD))=0$ for all sufficiently large $n$.

Fix a Picard component and a representative $L_0$ in it. The integral class $c_1(\OO_Y(S_j))$ is torsion because $H^2(Y,\Q)=0$; choose a positive integer $q$ killing it. Then $\OO_Y(-qnS_j)$ belongs to $\Pic^0(Y)$ for every $n$. Apply the preceding argument to $E=(\Omega_Y^1\otimes L_0)|_{S_i}$. It supplies, in the chosen Picard component, a twist $L_0\otimes\OO_Y(-qnS_j)$ with zero ambient group on $S_i$. Local holomorphic Picard families and proper semicontinuity show that the nonvanishing locus on that component is a proper closed analytic subset. Its complement is dense and open. The same argument works in every Picard component.

Choose an ordered list $A_1,\ldots,A_N$ of the prime divisors $S_i$, with $S_i$ repeated $m_i$ times, and put $F_0=0$, $F_k=A_1+\cdots+A_k$. The exact Cartier sequences, tensored with the locally free ambient bundle, are
\[
\begin{split}
0\longrightarrow
 (\Omega_Y^1\otimes L\otimes\OO_Y(-F_{k-1}))|_{A_k}
 \longrightarrow (\Omega_Y^1\otimes L)|_{F_k}\\
 \longrightarrow(\Omega_Y^1\otimes L)|_{F_{k-1}}
 \longrightarrow0.
\end{split}
\]
This follows locally from multiplication by an equation of $F_{k-1}$, and remains exact when a component occurs repeatedly. For each $k$ take the dense open set just proved on $A_k$ and translate it by the fixed global line $\OO_Y(-F_{k-1})$, using the appropriate Picard component. Their finite intersection is dense open. Every displayed first term has no sections there, so induction from $F_0=0$ proves (A.15b) on the full scheme $F_N=F$.

A proper holomorphic map from a smooth compact manifold to a smooth curve has finitely many critical values, since its proper analytic critical-value set does not contain a general point. Its smooth connected fibres are irreducible. Thus there are finitely many reducible fibres, and a further finite intersection gives the asserted simultaneous vanishing. This proves the final consequence and resolves the unequal-multiplicity transfer in this case without any assertion about torsion-free differentials on individual components.
\end{proof}

\subsection{The specified nonnormal fibre and the proof of Proposition 7.3}

\begin{statement}{A.4.1}{Proved}{Both fibre groups on the specified $W$}
For the reduced normal-crossings fibre $W$ with normalization $\eta:\operatorname{dP}_6\to W$ in the specified threefold $X$, and every $L\in\Pic(X)$, put $A=(K_X\otimes L^{-1})|_W$. Then
\[
H^0(W,Q\otimes A)=0,
\qquad H^0(W,\Omega_X^1|_W\otimes A)\simeq\C.
\tag{A.15}
\]
The same two conclusions hold with $A=B|_W$ for every $B\in\Pic(X)$.
\end{statement}
\begin{proof}
The only global topological input is $H^2(X,\Z)=0$, and the only identification of the fibre used is its specified normalization as the blowup of $\PP^2$ at three points. A blowup at a smooth point has $\pi_*\OO=\OO$ and $R^1\pi_*\OO=0$. The latter follows from the theorem on formal functions: successive exceptional-neighborhood quotients are $\OO_{\PP^1}(k)$, $k\ge0$, whose first cohomology is zero. Thus $H^1(\operatorname{dP}_6,\OO)=H^1(\PP^2,\OO)=0$. The exponential sequence makes $c_1:\Pic(\operatorname{dP}_6)\to H^2(\operatorname{dP}_6,\Z)$ injective. The pullback of every line bundle on $X$ has zero first Chern class and is therefore trivial.

The differential of normalization gives $\Omega_W^1\to\eta_*\Omega^1_{\operatorname{dP}_6}$. Its target is torsion-free as an $\OO_W$-module, so it factors through $Q$. The resulting map is injective because it is an isomorphism on the dense regular locus and its source is torsion-free. A holomorphic one-form on the blowup restricts away from the exceptional curves to a one-form on $\PP^2$ minus three points. Its coefficients extend over these points by Hartogs' theorem, and $H^0(\PP^2,\Omega^1_{\PP^2})=0$ by the Euler sequence. Hence $H^0(\operatorname{dP}_6,\Omega^1)=0$. The projection formula and the triviality of $\eta^*A$ prove the first equality in (A.15).

Because $W$ is the full reduced fibre, its conormal is trivial. A.1.2 identifies $K\otimes A$ with $\eta_*\eta^*A=\eta_*\OO_{\operatorname{dP}_6}$. The normalization is connected and compact, so its global holomorphic functions are $\C$. The cohomology sequence of A.1.1 and the first equality in (A.15) identify this group with the ambient group. The same argument applies to every $B$.
\end{proof}

\begin{statement}{A.4.2}{Refuted}{A specific implication in the printed proof of Proposition 7.3}
The implication that a ruled minimal surface $S_0$ with a one-parameter family of numerically trivial line bundles $M_t$ and nonzero sections of $\Omega^1_{S_0}\otimes M_t$ must be ruled over a curve of genus at least two is false when the family $M_t$ is allowed to be constant, as the printed proof explicitly allows.
\end{statement}
\begin{proof}
Take $S_0=E\times\PP^1$ for an elliptic curve $E$ and take $M_t=\OO_{S_0}$ for every parameter. The product cotangent decomposition gives
\[
 h^q(S_0,\Omega^1_{S_0})=(1,2,1),\qquad q=0,1,2.
\]
Indeed the $E$ cotangent summand contributes $(1,1,0)$, while the $\PP^1$ cotangent summand contributes $(0,1,1)$ by the Kunneth formula and $h^1(\PP^1,\OO(-2))=1$. A nonzero invariant form on $E$ provides the asserted sections for every parameter, whereas the ruling has genus one. This refutes the stated implication. It is a counterexample to that implication with its explicit data; it supplies no threefold satisfying all of Proposition~7.3's ambient hypotheses.
\end{proof}

\begin{statement}{A.5.1}{Proved}{What trivial normalization pullback forces}
Let $S$ be a connected reduced full fibre of a proper map from a complex manifold to a smooth curve. For every line bundle $A$ on $S$ with $\eta^*A\simeq\OO_{\widetilde S}$, the ambient group contains the canonical image of $H^0(\widetilde S,\OO_{\widetilde S})$. In particular it is nonzero. This applies to the subcase of Proposition~7.3 in which all $\eta^*(L_t|_S)$ are trivial. It is compatible with the vanishing of $H^0(S,Q\otimes A)$ through the exact maps of A.1.1, as A.4.1 demonstrates.
\end{statement}
\begin{proof}
A local coordinate of the base trivializes $N^*$. Tensoring the injection in (A.3) by $A$ and applying the projection formula gives
$\eta_*\eta^*A\hookrightarrow K\otimes A\hookrightarrow V\otimes A$.
Taking global sections preserves injectivity. Every connected component of the compact normalization has the nonzero constant function $1$, proving nonvanishing. No normal-crossings assumption is needed for this inclusion.
\end{proof}

The line in the printed proof asserting $H^2(\widetilde S,\omega_{\widetilde S})=0$ also cannot hold for a connected compact normal Cohen--Macaulay surface with its dualizing sheaf: Serre duality identifies its dual with $H^0(\widetilde S,\OO)=\C$. Replacing $\omega_{\widetilde S}$ there by $\OO_{\widetilde S}$ matches the vanishing stated at the beginning of that proof, but does not settle its omitted genus-one case. We have established (11) for $W$ and the exact obstruction to deducing ambient vanishing from it. The full statement of Proposition~7.3 is proved below in A.4.9 by a different argument through unitary local systems. The following argument establishes Proposition~4.11 independently of the monodromy-dependent restrictions.

\begin{statement}{A.5.2}{Proved}{Normal Gorenstein surfaces with vanishing first Chern class and Euler characteristic}
A compact connected normal Gorenstein complex surface $S$ with $c_1(\omega_S)=0$ in $H^2(S,\Q)$ and $\chi(S,\OO_S)=0$ is smooth. Consequently every singular irreducible reduced divisor in a compact complex threefold $Y$ with $H^2(Y,\Q)=0$ is nonnormal. In particular the assertion of CDP20 Proposition~4.11 holds without its monodromy-dependent restrictions on the smooth fibre.
\end{statement}
\begin{proof}
Let $\pi:Z\to S$ be the minimal resolution and set $\ell=h^0(S,R^1\pi_*\OO_Z)$. Normality gives $\pi_*\OO_Z=\OO_S$, and the fibres have dimension at most one, so Leray gives $\chi(Z,\OO_Z)=-\ell$. We first prove $\ell=0$. The classification table in \cite[Theorem~2.2]{Peters2021}, together with Noether's formula and invariance of $\chi(\OO)$ under blowing up, shows that a smooth compact surface of negative holomorphic Euler characteristic is a blown-up ruled surface over a smooth curve $B$ of genus $g\ge2$. Indeed the sums $c_1^2+c_2$ in every other row of that table are nonnegative. Thus, when $\ell>0$, there is a ruling $p:Z\to B$ and $\chi(Z,\OO_Z)=1-g$, giving $\ell=g-1$.

An exceptional irreducible curve $E$ of $\pi$ dominating $B$ would have $h^1(E,\OO_E)\ge g$: normalize $E$, apply Riemann--Hurwitz to its finite map to $B$, and use the normalization exact sequence on $E$. The exact sequence $0\to\mathcal I_E\to\OO_Z\to\OO_E\to0$, pushed forward by $\pi$, gives a surjection $R^1\pi_*\OO_Z\to R^1\pi_*\OO_E$ because $R^2\pi_*\mathcal I_E=0$. It follows that $\ell\ge h^1(E,\OO_E)\ge g$, a contradiction. Hence every exceptional curve is vertical for $p$. The ruling is constant on each connected fibre of $\pi$ and descends to a holomorphic map $p_S:S\to B$: local holomorphic coordinate functions descend by $\pi_*\OO_Z=\OO_S$, and properness supplies the required neighborhoods. The map $p$ is the composite of blowups and a $\PP^1$-bundle, so $R^1p_*\OO_Z=0$. The Leray sequence for $p=p_S\circ\pi$ then gives
\[
0=R^1p_*\OO_Z\longrightarrow (p_S)_*R^1\pi_*\OO_Z
 \longrightarrow R^2(p_S)_*\OO_S=0.
\]
A nonzero sheaf supported at finitely many points cannot have zero direct image, so again $\ell=0$, a contradiction. This proves rationality of all singularities.

We give the crepancy argument to specify the comparison of canonical bundles. Proper duality applied to $R\pi_*\OO_Z=\OO_S$ identifies $\pi_*\omega_Z$ with the invertible dualizing sheaf $\omega_S$. A local generator of $\omega_S$ therefore pulls back to a holomorphic two-form on $Z$, agreeing with the usual pullback off the exceptional set. Its zero divisor is effective and exceptional. These divisors glue to $D\ge0$ with $K_Z=\pi^*K_S+D$. Every exceptional irreducible curve has arithmetic genus zero by the same surjection to its first cohomology, hence is a smooth rational curve. Minimality of the resolution gives $E^2\le-2$, and adjunction gives $D\cdot E=K_Z\cdot E=-2-E^2\ge0$. The exceptional intersection matrix is negative definite. Thus a nonzero $D$ would satisfy both $D^2<0$ and $D^2=\sum_E d_E(D\cdot E)\ge0$. Therefore $D=0$.

Consequently $c_1(K_Z)=0$ in rational cohomology. It follows by adjunction that $Z$ contains no $(-1)$-curve, so $Z$ is a minimal compact surface; Noether's formula gives $c_2(Z)=0$. We exclude rational curves on each surviving row of the classification. A ruled surface is excluded by $K_Z\cdot F=-2$ on a ruling fibre. General-type surfaces have $K_Z^2>0$ and K3 and Enriques surfaces have positive $c_2$, so none occurs. A torus has no rational curve because a map from $\PP^1$ lifts to its vector-space universal cover and is constant; the same holds for a finite etale quotient of a torus. A primary Kodaira surface is an elliptic bundle over an elliptic curve, so a rational curve would map constantly to the base and then constantly to a fibre; finite etale quotients have the same property. A class VII surface here has $b_2=c_2=0$, and the Hopf--Inoue classification gives no rational curves. Equivalently one can use their standard universal covers: $\C^2\setminus\{0\}$ for a primary Hopf cover and $\mathbb H\times\C$ for an Inoue surface. A map from $\PP^1$ lifts and has constant coordinate functions.

It remains to check the properly elliptic row. For a minimal elliptic surface, the Kodaira fibre list expresses $c_2$ as a sum of nonnegative fibre Euler numbers, with zero contributions exactly at multiples of smooth elliptic curves. Since $c_2(Z)=0$, all fibre reductions are smooth elliptic curves. The canonical bundle formula on a Kähler elliptic surface with $\chi(\OO_Z)=0$ has rational first Chern class
\[
 c_1(K_Z)=\left(2g(B)-2+\sum_i(1-m_i^{-1})\right)[F].
\]
In Kodaira dimension one the coefficient is positive, and the Kähler class has positive integral over $F$, contradicting $c_1(K_Z)=0$. In the non-Kähler case, a horizontal rational curve would be a multisection $C$. For $D=C+nF$ with $n$ sufficiently large, $D^2>0$ and $(K_Z-mD)\cdot F<0$ for $m>0$. Riemann--Roch therefore gives quadratic growth of $h^0(Z,\OO_Z(mD))$, making $Z$ Moishezon and hence projective as a smooth compact surface. This contradicts the non-Kähler case. Vertical rational curves have already been excluded by the fibre list. Thus $Z$ has no rational curves in any case. The exceptional curves of $\pi$ were rational, so there are none and $S$ is smooth.

Finally take an irreducible reduced divisor $S\subset Y$ with $H^2(Y,\Q)=0$. It is a Cartier divisor in the smooth manifold $Y$, hence Gorenstein, and adjunction gives $\omega_S=(K_Y\otimes\OO_Y(S))|_S$, with rational first Chern class zero. The divisor sequence and Hirzebruch--Riemann--Roch give
$\chi(\OO_S)=\chi(\OO_Y)-\chi(\OO_Y(-S))=0$, since $c_1(\OO_Y(S))=0$ rationally. A normal $S$ is therefore smooth by the result just proved.
\end{proof}

\begin{statement}{A.5.3}{Proved}{The effective conductor divisor on a minimal resolution}
Let $S$ be a compact reduced hypersurface in a smooth threefold, let $\eta:\coprod_i T_i\to S$ be its normalization, and let $\pi_i:Z_i\to T_i$ be the minimal resolutions. Put $h_i=\eta\pi_i$. There are canonical effective integral divisors $D_i$ on $Z_i$ with
\[
 \omega_{Z_i}(D_i)\simeq h_i^*\omega_S.
 \tag{A.25}
\]
The nonexceptional part of $D_i$ is exactly the divisorial conductor, with its scheme-theoretic multiplicities. When $S\subset Y$ and $H^2(Y,\Q)=0$, every $K_{Z_i}+D_i$ has rational first Chern class zero. A ruled $Z_i\to B$ with $g(B)\ge2$ then has
\[
 D_i=2C+V,
 \tag{A.26}
\]
where $C$ is a section and $V$ is vertical. In the case where $C$ is $\pi_i$-exceptional, $\eta^*(L|_S)|_{T_i}$ is trivial for every $L\in\Pic(Y)$.
\end{statement}
\begin{proof}
Finite duality gives $\omega_{T_i}=\mathfrak c_i\otimes\eta^*\omega_S$ as rank-one reflexive sheaves, where $\mathfrak c_i$ is the conductor viewed as an ideal on $T_i$. Locally this is the equality
$\Hom_R(B,R)=\{q:qB\subset R\}=\mathfrak c$ for a reduced Gorenstein ring $R$ and its finite normalization $B$, tensored with the invertible dualizing module. The equality uses the meromorphic multiplication identification as in A.2.1. The conductor is a maximal Cohen--Macaulay module because it is this dual of the maximal Cohen--Macaulay normalization over the Gorenstein ring. On a normal surface a rank-one maximal Cohen--Macaulay module is reflexive, so the displayed equality is determined by its codimension-one orders. Those orders define an effective conductor divisor $D_i^0$ on $T_i$.

Off the exceptional set the canonical identification pulls back to $\omega_{Z_i}(D_i^{\rm str})\simeq h_i^*\omega_S$, where $D_i^{\rm str}$ denotes strict transform. The meromorphic extension of this line-bundle identification gives integral exceptional coefficients $a_j$ with
$K_{Z_i}+D_i^{\rm str}+\sum_j a_jE_j=h_i^*K_S$.
We verify their signs. For each exceptional irreducible curve $E_j$, minimality excludes a smooth rational $(-1)$-curve. Adjunction therefore gives
$K_{Z_i}\cdot E_j=2p_a(E_j)-2-E_j^2\ge0$;
also $D_i^{\rm str}\cdot E_j\ge0$. The exceptional intersection matrix $M=(E_j\cdot E_k)$ is negative definite with nonnegative off-diagonal entries, and
\[
 M a=-v,\qquad v_j=(K_{Z_i}+D_i^{\rm str})\cdot E_j\ge0.
\]
To see that $a\ge0$ without an unproved sign convention, write $a=a_+-a_-$ with disjoint nonnegative coordinate supports. If $a_-\ne0$, then
$a_-^tMa=a_-^tMa_+-a_-^tMa_->0$,
whereas $Ma=-v$ gives $a_-^tMa\le0$. This contradiction proves $a_j\ge0$. Thus $D_i=D_i^{\rm str}+\sum_j a_jE_j$ proves (A.25). The conductor duality and effective exceptional comparison also appear, for projective surfaces with trivial dualizing sheaf, in \cite[Lemma~1.1 and Corollary~1.4]{Furushima1993}; the argument here supplies the maps and does not use that paper's additional K3-type hypothesis $h^1(\OO_S)=0$.

For a divisor in $Y$, adjunction gives $\omega_S=(K_Y\otimes\OO_Y(S))|_S$, proving the numerical assertion. On a general ruling fibre $f$, $K_{Z_i}\cdot f=-2$, so $D_i\cdot f=2$. Every horizontal irreducible curve has normalization of genus at least $g(B)\ge2$. Such a curve $C_0$ cannot occur in $D_i$ with coefficient one: otherwise
\[
 0=(K_{Z_i}+D_i)\cdot C_0
   =2p_a(C_0)-2+(D_i-C_0)\cdot C_0>0.
\]
Every horizontal coefficient is consequently at least two. The degree-two intersection with $f$ forces exactly one horizontal component, coefficient two and degree one over $B$, giving (A.26). The degree-one finite morphism $C\to B$ is an isomorphism because $B$ is normal; thus $C$ is a smooth section.

Finally, integral $H^2$ of a smooth ruled surface and of its blowups is torsion-free, and its topologically trivial line bundles are precisely the pullbacks of $\Pic^0(B)$. These facts follow from the projective-bundle and blowup cohomology formulas and from $p_*\OO_{Z_i}=\OO_B$, $R^1p_*\OO_{Z_i}=0$. Restriction to the section $C$ is inverse to this Picard pullback. The line $h_i^*L$ has torsion, hence zero, integral first Chern class. When $C$ is exceptional, its restriction to $C$ is trivial because $h_i$ maps $C$ to a point. Therefore $h_i^*L$ is trivial on $Z_i$. Normality gives $(\pi_i)_*\OO_{Z_i}=\OO_{T_i}$, and the projection formula descends this triviality to $T_i$. This proves the last assertion.
\end{proof}

\begin{statement}{A.5.4}{Proved}{The remaining normalization class for a persistent ambient obstruction}
Let $Y$ satisfy the CDP hypotheses, and let $S$ be a connected nonnormal reduced fibre support for which
\[
 H^0(S,\Omega_Y^1|_S\otimes L|_S)\ne0
 \quad\hbox{for every }L\in\Pic(Y).
 \tag{A.27}
\]
There is a normalization component $T_i$ whose minimal resolution $Z_i$ has Kodaira dimension $-\infty$. Every line from $\Pic(Y)$ pulls back to a unitary-flat line on that $Z_i$. A rational or class VII $Z_i$ has trivial Picard pullback from $Y$, and then
\[
 H^2\!\left(T_i,\omega_{T_i}\otimes
             \eta^*((K_Y^{-1}\otimes L)|_S)|_{T_i}\right)\simeq\C
 \quad\hbox{for every }L\in\Pic(Y).
 \tag{A.28}
\]
The same conclusion holds for the exceptional-section case of A.5.3. Thus the only normalization type left by this argument, when (A.28) has not been obtained, is an irrational ruled surface with nonconstant Picard pullback; in base genus at least two its divisor (A.26) has its horizontal section in the strict conductor.
\end{statement}
\begin{proof}
A.4.9 gives $H^0(S,\widetilde\Omega_S^1\otimes L|_S)=0$ on a dense open subset of $\Pic(Y)$. By A.1.1 and (A.27), $H^0(S,K\otimes L|_S)$ is nonzero there. Its nonvanishing locus is closed analytic in the holomorphic Picard family, hence is all of $\Pic(Y)$.

Pull $K$ to each $Z_i$, quotient the resulting sheaf by its torsion, and take its double dual. The result is a fixed holomorphic line bundle $M_i$, since its generic rank on each component is one and $Z_i$ is a smooth surface. The induced maps give an injection
\[
 H^0(S,K\otimes L|_S)\hookrightarrow
       \bigoplus_i H^0(Z_i,M_i\otimes h_i^*L).
 \tag{A.29}
\]
Indeed a section in the kernel vanishes generically on every component, and $K\subset\Omega_Y^1|_S$ is torsion-free over the reduced space. Each summand's nonvanishing locus is closed analytic in the connected complex vector space $\Pic(Y)=H^1(Y,\OO_Y)$. A finite union of proper analytic subsets cannot cover that space. Consequently one component $i$ has
$H^0(Z_i,M_i\otimes h_i^*L)\ne0$ for every $L$.

Fix a Gauduchon metric on $Z_i$, with fundamental form $\omega$ satisfying $\partial\bar\partial\omega=0$. The degree of a line bundle is the integral of its real Chern curvature against $\omega$; it is independent of the chosen Hermitian metric and additive under tensor product. A nonzero section has an effective zero divisor, so its line has nonnegative degree. For any $G\in\Pic(Y)$ the preceding nonvanishing applied to $G^{\otimes n}$ therefore gives
\[
 \deg_\omega(M_i)+n\deg_\omega(h_i^*G)\ge0
 \quad\hbox{for every }n\in\Z.
\]
It follows that $\deg_\omega(h_i^*G)=0$ for every $G$ and every Gauduchon metric.

We record the precise flatness implication. For a topologically trivial line, the real Chern curvature $F$ is exact. If $\int F\wedge\omega=0$, the Gauduchon scalar elliptic equation supplies a real $u$ with $\Lambda_\omega(F-i\partial\bar\partial u)=0$. The remaining form is closed, exact, real and primitive of type $(1,1)$, hence anti-selfdual. Its square integrates both to zero by exactness and to minus its squared norm. It is therefore zero. Changing the line metric by $e^{-u}$ gives zero Chern curvature, hence a flat unitary connection. This is exactly the curvature argument of \cite[Lemma~2.3(1)--(2)]{Teleman2006}. Since $H^2(Y,\Z)=0$, it applies to every $h_i^*G$.

The divisor $D_i$ of A.5.3 is nonzero. For an irreducible nonnormal support this follows from its nonempty conductor; for a connected reducible support every component meets another component, again giving a nonempty conductor. If $Z_i$ had nonnegative Kodaira dimension, some positive pluricanonical power would have a nonzero section, making $\deg_\omega(K_{Z_i})\ge0$. But $D_i$ is a nonzero effective divisor and therefore has strictly positive degree. Equation (A.25), with $\omega_S=(K_Y\otimes\OO_Y(S))|_S$, would give
\[
 0=\deg_\omega h_i^*(K_Y\otimes\OO_Y(S))
   =\deg_\omega(K_{Z_i})+\deg_\omega(D_i)>0,
\]
a contradiction. Thus $\kappa(Z_i)=-\infty$, whose classification consists of rational, irrational ruled, and class VII surfaces.

For a rational $Z_i$, topologically trivial lines are trivial. For a class VII $Z_i$, $h^1(\OO_{Z_i})=1$ and $b_1(Z_i)=1$. With our exponential convention $u\mapsto\exp(2\pi\mathrm i u)$, the unitary-flat subgroup of $\Pic^0(Z_i)$ has real tangent space equal to the injective image of $H^1(Z_i,\mathbb R)$ in $H^1(Z_i,\OO_{Z_i})$. This is \cite[Lemma~2.3(2)]{Teleman2006} after multiplication by the fixed scalar $\mathrm i$, or follows from the pluriharmonic maximum principle. That real tangent space has dimension one. The holomorphic Picard pullback has complex-linear differential and, by the proved flatness, takes values in this real one-dimensional subspace. Its differential is zero, so the homomorphism is constant and equals zero at the trivial line. Every $h_i^*L$ is therefore trivial. Projection formula and $(\pi_i)_*\OO_{Z_i}=\OO_{T_i}$ again descend triviality to $T_i$. The exceptional-section case has the same conclusion by A.5.3.

In each of these cases, the line multiplying $\omega_{T_i}$ in (A.28) is trivial, since it too comes from $Y$. A normal complex surface is Cohen--Macaulay, so its dualizing sheaf and Serre duality give $H^2(T_i,\omega_{T_i})=H^0(T_i,\OO_{T_i})^\vee=\C$. This proves (A.28) and the stated exact reduction of the remaining normalization problem.
\end{proof}

\begin{statement}{A.5.6}{Proved}{Uniform negative-degree vanishing for a fixed vector bundle}
Let $Z$ be a compact smooth complex surface, $\omega$ a Gauduchon metric, and $E$ a fixed holomorphic vector bundle. There is a real constant $b(E,\omega)$ such that
\[
 \deg_\omega P<b(E,\omega)
 \quad\Longrightarrow\quad H^0(Z,E\otimes P)=0
 \tag{A.30}
\]
for every holomorphic line bundle $P$. Thus persistence of sections of $E\otimes h^*L$ for every $L\in\Pic(Y)$ forces $\deg_\omega h^*L=0$ for every $L$, for any holomorphic map $h:Z\to Y$. Persistence on one Picard component already forces degree zero on the pullback of $\Pic^0(Y)$. Under $H^2(Y,\Q)=0$, it forces degree zero and unitary flatness for every ambient Picard pullback. On a class VII surface $Z$ with $H^2(Y,\Z)=0$, the entire holomorphic Picard pullback is the zero map.
\end{statement}
\begin{proof}
Choose a Hermitian metric on $E$ and let $C_E$ bound the largest eigenvalue of the Hermitian curvature endomorphism $\mathrm i\Lambda_\omega F_E$ on the compact surface. Write $V_\omega=\int_Z\omega^2/2$ and normalize degree by
$\deg_\omega P=(2\pi)^{-1}\int_Z\mathrm iF_P\wedge\omega$.
The scalar Gauduchon elliptic equation changes a Hermitian line metric so that its contracted curvature is the constant
\[
 \mathrm i\Lambda_\omega F_P
       =c_P:=2\pi\deg_\omega(P)/V_\omega.
\]
Indeed the difference from this constant has integral zero, and the image of $u\mapsto\mathrm i\Lambda_\omega\partial\bar\partial u$ is precisely the mean-zero functions for a Gauduchon metric; this is the scalar equation used in \cite[Lemma~2.3]{Teleman2006}. The Hermitian line metric may therefore be rescaled by the corresponding $e^{-u}$.

For a holomorphic section $s$ of $E\otimes P$, the Chern-connection identity is
\[
 \mathrm i\Lambda_\omega\partial\bar\partial|s|^2
 =|\nabla^{1,0}s|^2
  -\langle\mathrm i\Lambda_\omega F_{E\otimes P}s,s\rangle.
\]
It follows directly in a holomorphic frame by differentiating the Hermitian norm; thus no Kähler identity is being assumed. At a positive maximum of $|s|^2$, the left side is nonpositive, whereas the right side is at least
$-(C_E+c_P)|s|^2$.
If $c_P<-C_E$, this is strictly positive, a contradiction. We may take $b(E,\omega)=-C_EV_\omega/(2\pi)$, proving (A.30).

Apply this bound to $P=h^*(L^{\otimes n})$ for every positive and negative integer $n$. Nonzero sections for all $n$ are incompatible with a nonzero degree, proving the persistence assertion. For persistence on a single component, absorb one representative line into the fixed bundle $E$ and apply the same argument to all powers of every line in $\Pic^0(Y)$. If $H^2(Y,\Q)=0$, a positive power of every ambient line lies in $\Pic^0(Y)$, so additivity of degree gives degree zero for every pullback. Its real first Chern class is zero, even if its integral class is torsion. The Chern curvature is therefore exact, and the exact-curvature argument in A.5.4 again produces a flat unitary connection. On a class VII surface the differential of the holomorphic Picard pullback then lies in a real one-dimensional unitary tangent space. Complex linearity makes that differential zero. The exponential sequence makes $\Pic(Y)$ a connected quotient of $H^1(Y,\OO_Y)$ when $H^2(Y,\Z)=0$, so under that integral hypothesis the entire map is identically zero.
\end{proof}

\begin{statement}{A.5.9}{Proved}{The ruled obstruction with a reduced conductor image}
Let $S\subset Y$ be an irreducible nonnormal reduced Cartier
divisor in a compact smooth complex threefold with
$H^2(Y,\Z)=0$ and $\chi(\OO_S)=0$.
Let $\eta:T\to S$ be its normalization and
$\pi:Z\to T$ its minimal resolution, and put $h=\eta\pi$.
Suppose that $Z$ is irrational ruled, with ruling $p:Z\to B$,
that every exceptional curve of $\pi$ is vertical, and that
the effective divisor of A.5.3 has the form
\[
 K_Z+D_Z=h^*\omega_S,\qquad D_Z=2C+V,                 \tag{A.RC.1}
\]
where $C$ is a section in the strict conductor and $V$ is
vertical. Suppose also that the conductor image scheme
$\Gamma=\operatorname{Spec}_S(\OO_S/\mathfrak c)$ is reduced,
and that $H^0(S,K\otimes L|_S)\ne0$ for every $L\in\Pic(Y)$,
with $K$ as in A.1.1.
Then every ambient line bundle pulls back trivially to $T$.
In particular
\[
 H^2\bigl(T,\omega_T\otimes
       \eta^*((K_Y^{-1}\otimes L)|_S)\bigr)\simeq\C
       \quad\hbox{for every }L\in\Pic(Y).            \tag{A.RC.2}
\]
The proof forces $B$ and $\Gamma$ to be smooth elliptic curves,
$Z=T$ to be minimal ruled, and $D_Z=2C$ with $C^2=0$.
Thus no local-complete-intersection assumption on the original
conductor image curve is needed in this reduced-scheme case.
\end{statement}
\begin{proof}
Write $\mathcal B=\eta_*\OO_T$ and
$\mathcal Q=\mathcal B/\OO_S$. We first specify conductor
duality, including its twists. Since $S$ is Gorenstein of
dimension two and $\mathcal B$ is maximal Cohen--Macaulay,
its dual is $\mathfrak c=\mathcal Hom_S(\mathcal B,\OO_S)$,
and maximal Cohen--Macaulay biduality gives
$\mathcal Hom_S(\mathfrak c,\OO_S)=\mathcal B$.
The quotient $\OO_\Gamma$ is Cohen--Macaulay of dimension one:
this follows from the depth lemma in
$0\to\mathfrak c\to\OO_S\to\OO_\Gamma\to0$.
Applying $\mathcal Hom_S(-,\OO_S)$ to that sequence, its
initial terms identify the actual normalization inclusion and
its cokernel as
\[
 0\longrightarrow\OO_S\longrightarrow\mathcal B
 \longrightarrow
 \mathcal Ext^1_S(\OO_\Gamma,\OO_S)\longrightarrow0,
 \qquad
 \mathcal Q\simeq\omega_\Gamma\otimes\omega_S^{-1}|_\Gamma.
                                                               \tag{A.RC.3}
\]
Adjunction makes $\omega_S$ the restriction of
$K_Y\otimes\OO_Y(S)$. Its degree is zero on every compact
curve, because $H^2(Y,\Z)=0$. Curve Riemann--Roch and Serre
duality therefore give
$\chi(\mathcal Q)=\chi(\omega_\Gamma)=-\chi(\OO_\Gamma)$.
On the other hand, A.5.8 gives rational singularities of $T$
and $\chi(\OO_T)=1-g(B)$. The normalization sequence and
$\chi(\OO_S)=0$ now give the exact numerical identity
\[
                   \chi(\OO_\Gamma)=g(B)-1.          \tag{A.RC.4}
\]

We next prove that $\Gamma$ is connected. Contract ruling
$(-1)$-curves from $Z$ to a minimal ruled bundle $Z_0\to B$.
Under a single contraction $\sigma:Z_1\to Z_2$ with exceptional
curve $E$, let $D_2=\sigma_*D_1$, let $a$ be the coefficient
of $E$ in $D_1$, and let $m$ be the multiplicity of $D_2$ at
the contraction point. The canonical divisor formula gives
\[
 K_{Z_1}+D_1=\sigma^*(K_{Z_2}+D_2)+(1+a-m)E.
                                                               \tag{A.RC.5}
\]
Numerical triviality and intersection with $E$ force $a=m-1$.
The same numerical triviality descends to $Z_2$. On $Z_0$
the pushed divisor is $2C_0+V_0$, with $C_0$ a section and
$V_0$ a nonnegative sum of fibres, so its support is connected.
Reverse the contractions. A centre outside the existing support
would produce coefficient $-1$ on its exceptional curve; that
coefficient persists on strict transforms and contradicts the
effectiveness of the final $D_Z$. A centre of multiplicity one
is a smooth point of one coefficient-one component and its
exceptional coefficient is zero, so it does not disconnect the
support. At a centre of multiplicity at least two the exceptional
curve has positive coefficient and connects all strict-transform
branches through that point. Induction proves that
$\operatorname{Supp}D_Z$ is connected.

Every positive exceptional component of $D_Z$ maps into
$\Gamma$. Indeed away from $\Gamma$ the surface $T$ equals
the Gorenstein surface $S$ and has rational singularities by
A.5.8. Such a singularity is crepant on its minimal resolution:
a generator of its dualizing sheaf pulls back holomorphically
by rational duality; its effective exceptional zero divisor
has nonnegative intersection with each exceptional curve by
adjunction and minimality, so negative definiteness makes it
zero, as in A.5.2. Thus (A.RC.1) has no positive exceptional
coefficient there. The strict conductor maps onto the whole
conductor support. It follows that
$h(\operatorname{Supp}D_Z)=\operatorname{Supp}\Gamma$,
which proves connectedness.

Let $D\subset\Gamma$ be the image of $C$. Localize (A.RC.3)
at the generic point of $D$. Because $\Gamma$ is reduced,
$\OO_\Gamma$ there is the function field $k(D)$, and its
dualizing module has dimension one over that field. The
elementary conductor quotient sequence is
\[
 0\longrightarrow\OO_\Gamma\longrightarrow
          \mathcal B/\mathfrak c
   \longrightarrow\mathcal Q\longrightarrow0.
\]
It follows that $\mathcal B/\mathfrak c$ has dimension two
over $k(D)$. The conductor coefficient two on $C$ contributes
$2\deg(C\to D)$ to this dimension. Therefore this degree is
one and there is no other normalization branch over $D$.
The generic transverse ring has conductor $t^2\C\{t\}$ and
is exactly $\C+t^2\C\{t\}=\C\{t^2,t^3\}$, the cusp ring.
Thus $D^\nu\simeq C\simeq B$. This deduction also rules out
any additional vertical branch over the same image curve.

A connected compact reduced curve has only constant global
holomorphic functions. From (A.RC.4) we therefore obtain
$h^1(\Gamma,\OO_\Gamma)=2-g(B)$.
Restriction to the reduced component $D$ gives a surjection
on $H^1$, because its kernel is a coherent sheaf on a curve
and has no $H^2$. Normalization of $D$ gives
$h^1(D,\OO_D)\ge g(B)$. Hence
\[
             2-g(B)\ge h^1(D,\OO_D)\ge g(B).         \tag{A.RC.6}
\]
As $g(B)\ge1$, all quantities equal one. The normalization
quotient of $D$ has length zero, so $D$ itself is smooth
elliptic. Statement C2d applies to this actual smooth curve
in $Y$ and makes the transversal cusp constant along it.
No other singular curve meets $D$. Connectedness of $\Gamma$
therefore gives $\Gamma=D$. Its normalization inverse image
is the single smooth curve $C$, the normalization is smooth
near it, and C2d gives $C^2=0$.

Adjunction in (A.RC.1) yields $V\cdot C=0$. A ruling
$(-1)$-curve outside $V$ cannot meet $C$, since its
intersection with $K_Z+2C+V$ would be at least $-1+2$.
A component of $V$ cannot meet $C$ either. Thus all
contractions to $Z_0$ avoid $C$ and preserve its square zero.
On the minimal bundle, $K_{Z_0}+2C_0$ is numerically trivial,
so its nonnegative vertical divisor $V_0$ is zero.
The first reverse blowup would be away from $C_0$ and have
coefficient $-1$ by (A.RC.5). There are no such blowups.
Thus $Z$ is minimal ruled and $D_Z=2C$.
An exceptional curve for $Z\to T$ would have negative square;
adjunction with $K_Z\equiv-2C$ would make it rational and
therefore vertical. A minimal ruled bundle has only fibres
of square zero, so no such exceptional curve exists and $Z=T$.

The constant-cusp computation A.2.4 now gives the injection
$K\hookrightarrow\eta_*(h^*\OO_Y(-S)\otimes\OO_Z(C))$.
Every ambient line is topologically trivial on $Z$ and is
the pullback of a degree-zero line on the elliptic curve $B$.
The section exact sequence, together with $R^1p_*\OO_Z=0$,
gives
\[
 0\longrightarrow\OO_B\longrightarrow p_*\OO_Z(C)
 \longrightarrow N_{C/Z}\longrightarrow0,
 \qquad\deg N_{C/Z}=0.                              \tag{A.RC.7}
\]
After a degree-zero twist $P$, both end terms have zero
sections unless $P=\OO_B$ or $P=N_{C/Z}^{-1}$.
The image of $\Pic(Y)$ in $\Pic^0(B)$ is connected:
$H^2(Y,\Z)=0$ makes $\Pic(Y)$ the connected image of
$H^1(Y,\OO_Y)$ under the exponential map.
All-twist persistence of $K$, and translation by the ambient
line $\OO_Y(-S)$, put this entire image in those at most two
points. It is therefore the identity. This proves triviality
of every ambient pullback. Serre duality on $T$ proves
(A.RC.2).
\end{proof}

\begin{statement}{A.5.11}{Proved}{The genus of a horizontal conductor image from a local torus group}
Let $f:Y\to C_0$ be a proper holomorphic map from a smooth
compact complex threefold to a smooth curve, with connected
fibres and complex two-torus general fibre. Let $S$ be an
irreducible reduced fibre support, $\eta:T\to S$ its
normalization, and $\pi:Z\to T$ a resolution which is
irrational ruled, with ruling $p:Z\to B$.
Suppose every exceptional curve of $\pi$ is vertical, and
the strict conductor has just one horizontal irreducible
component $C$, which is a section of $p$.
Let $D=\eta\pi(C)$ and let $\Sigma=D^\nu$ be its normalization.
Then $g(\Sigma)\le1$.
In particular, when $C\to D$ has degree one, $B$ is elliptic.
This applies to the horizontal-cusp case of $D_Z=2C+V$;
it allows arbitrary multiplicity of the full fibre and
arbitrary singularities or nonreduced conductor image on
the finitely many vertical curves.
\end{statement}
\begin{proof}
There is nothing to prove when $\Sigma$ is rational, so
assume $g(\Sigma)\ge1$.
Statement A.5.8 gives a holomorphic ruling
$\widetilde p:T\to B$ with $p=\widetilde p\pi$.
Identifying $C$ with $B$ by its section map, normalization
of its nonconstant map to $D$ gives a finite surjective
holomorphic map $\phi:B\to\Sigma$.
Put $\psi=\phi\widetilde p:T\to\Sigma$.

We describe precisely the finite identifications needed to
make this map descend to $S$. Every vertical conductor
component of $T$ lies over one of finitely many points of
$B$, so $\psi$ is constant on it. Its image under $\eta$
has rational normalization: a vertical component is the
image of a rational curve on $Z$, and a nonconstant map
from that curve to its image normalization is finite.
Such an image cannot equal $D$, because $g(\Sigma)\ge1$.
Consequently these vertical images meet $D$ at finitely
many points. Away from those points and from the finitely
many singular points of $D$, every normalization preimage
of a point of $D$ lies on $C$, and its $\psi$ value is the
unique corresponding point of $\Sigma$.
At a point outside $D$ with more than one normalization
preimage, all those preimages lie on vertical conductor
components. There are no other horizontal conductor
components by hypothesis. Thus there is a finite set
$A\subset\Sigma$ such that, for every $s\in S$, either
$\psi(\eta^{-1}(s))$ is a singleton or it is contained
in $A$. One can take $A$ to contain the constant values
on all vertical conductor components and all points of
$\Sigma$ above the singularities of $D$ and its
intersections with the vertical image curves.

Choose an embedded finite tree $J$ in the underlying
oriented surface $\Sigma$ containing $A$, and let
$q:\Sigma\to\Sigma/J$ collapse that tree to one point.
Such a tree is obtained by joining the finite set by
arcs and retaining a spanning tree after subdividing
their finitely many crossings. Triangulate the surface
with $J$ a subcomplex. A contractible subcomplex is a
cofibration; contracting it gives a homotopy equivalence
of this CW pair. Concretely the homotopy extension
property extends a contraction of $J$ and supplies a
homotopy inverse to the quotient map. In particular
\[
              q_*:\pi_1(\Sigma)\xrightarrow{\sim}
                         \pi_1(\Sigma/J).          \tag{A.Pi.1}
\]
The map $q\psi$ is constant on every fibre of the finite
proper surjection $\eta$. A proper continuous surjection
between these compact Hausdorff spaces is a quotient
map, so there is a unique continuous map
\[
 a:S\longrightarrow\Sigma/J,
 \qquad a\eta=q\psi.                              \tag{A.Pi.2}
\]
Let $\delta:\Sigma\to D\subset S$ be the curve normalization
map. On its dense smooth locus, lifting through $C$ gives
$a\delta=q$. Continuity gives this equality everywhere.
Equation (A.Pi.1) therefore makes $a_*$ surjective on
fundamental groups, without any assertion that a finite
point-identification quotient has the same fundamental
group as $\Sigma$.

We next extend the actual map (A.Pi.2) to a neighbourhood
of $S$ in $Y$. The simultaneous analytic triangulation
theorem of \cite[\S3, Theorem~2, p.~464]{Lojasiewicz1964}
triangulates the pair $(Y,S)$, since a complex analytic
subset is real semianalytic. Compactness makes the
subcomplex representing $S$ finite. After barycentric
subdivision it is full: vertices lying in it span a
face in that same subcomplex. Write $\lambda_v$ for
barycentric coordinates and let $L$ be its vertex set.
On the open neighbourhood where $\sum_{v\in L}\lambda_v>0$,
the formula
\[
 r(x)=\frac{\sum_{v\in L}\lambda_v(x)v}
                  {\sum_{v\in L}\lambda_v(x)}       \tag{A.Pi.3}
\]
defines a continuous retraction to $S$ in triangulation
coordinates. Its value lies in the face spanned by the
$L$ vertices of the same simplex, so these formulas
agree across faces. This supplies an actual open
neighbourhood $N$ and a map $r:N\to S$ fixed on $S$.

Let $c=f(S)$. Properness allows a small disc $\Delta$
about $c$ with $U=f^{-1}(\Delta)\subset N$ and no other
critical value. The map $(ar)|_U:U\to\Sigma/J$ is
surjective on fundamental groups: its restriction to
$\delta(\Sigma)\subset S\subset U$ is $q$, whose
fundamental-group map is an isomorphism.
Set $U^*=U\setminus S=f^{-1}(\Delta^*)$.
The inclusion induces a surjection
$\pi_1(U^*)\twoheadrightarrow\pi_1(U)$.
Indeed a loop in the smooth six-manifold $U$ can be
perturbed relative to its base point to be transverse
to an analytic stratification of $S$; every stratum
has real codimension at least two, so the resulting
one-dimensional loop misses every stratum.

Over $\Delta^*$, the proper submersion is a smooth
torus bundle. The fundamental-group exact sequence
and $\pi_2(\Delta^*)=0$ give
\[
 1\longrightarrow\Z^4\longrightarrow\pi_1(U^*)
       \longrightarrow\Z\longrightarrow1.
\]
A lift of the generator of $\Z$ splits this group
extension, hence
\[
 \pi_1(U^*)\simeq\Z^4\rtimes_T\Z
 \twoheadrightarrow\pi_1(U)
 \twoheadrightarrow\pi_1(\Sigma/J)
 \simeq\pi_1(\Sigma).                             \tag{A.Pi.4}
\]
The first group is metabelian: its commutator subgroup
is contained in the abelian normal subgroup $\Z^4$.
Every quotient in (A.Pi.4) is therefore metabelian.
A closed orientable surface of genus at least two
has a fundamental group surjecting onto the free
group on two generators: in its usual surface
presentation, retain two $a$ generators and send all
$b$ generators and the other $a$ generators to the
identity. That free group surjects onto $S_4$ by
sending its generators to $(1234)$ and $(12)$;
their conjugates supply the adjacent transpositions.
The commutator subgroup of $S_4$ is the nonabelian
group $A_4$, so this quotient is not metabelian.
This contradicts (A.Pi.4). Thus $g(\Sigma)=1$ in the
positive-genus case. When $C\to D$ has degree one,
$\Sigma\simeq C\simeq B$, which proves the final
assertion.
\end{proof}

\begin{statement}{A.5.12}{Proved}{The arithmetic genus of conductor images from coherent duality}
Let $f:Y\to C_0$ be a proper connected-fibre map from a
smooth compact complex threefold with $H^2(Y,\Z)=0$ to
a smooth curve, and let its smooth general fibre have
trivial canonical bundle. This includes complex
two-tori and primary Kodaira surfaces.
Let $mS$ be a full fibre with irreducible
reduced nonnormal support $S$. Write $\eta:T\to S$ for
normalization and $\pi:Z\to T$ for a resolution.
If the pullback homomorphism
$\Pic(Y)\to\Pic(Z)$ is nonconstant, then every reduced
irreducible component $D$ of the conductor image scheme
has arithmetic genus at most one.
In particular every such component whose normalization
has positive genus is a smooth elliptic curve.
The conductor image scheme itself is allowed to be
nonreduced, and no local-complete-intersection assumption
is imposed on its reduced components.
\end{statement}
\begin{proof}
First we determine exactly the needed canonical line on
the reduced support, without a multiplicity formula for
the direct image. The sheaf $f_*K_Y$ is torsion-free of
generic rank one on the smooth curve and hence is a
line bundle. Choose a frame near $c=f(S)$ and evaluate
it in $K_Y$. A coordinate $t$ on the base disc supplies
the conormal frame $dt$, and adjunction identifies
$K_Y|_F$ with $K_F$ after dividing by this frame.
Under that line-bundle isomorphism, the evaluated
section restricts on a general smooth fibre to a
nonzero canonical two-form, which is nowhere zero.
Its zero divisor is therefore
vertical. Shrinking the base disc to remove the other
vertical zero divisors, irreducibility of the support
gives an effective divisor $kS$, with $k\ge0$, and an
actual line-bundle comparison on $U=f^{-1}(\Delta)$:
\[
 K_Y|_U\simeq\OO_U(kS),\qquad
 \omega_S\simeq\OO_Y((k+1)S)|_S,\qquad
 \omega_S^{\otimes m}\simeq\OO_S.                 \tag{A.CD.1}
\]
The last isomorphism uses the local base equation whose
zero divisor is exactly the full fibre $mS$.
In particular no assertion $k=m-1$ is used.

Hirzebruch--Riemann--Roch in
$0\to\OO_Y(-S)\to\OO_Y\to\OO_S\to0$ gives
$\chi(\OO_S)=0$: the rational Chern character of
$\OO_Y(-S)$ equals that of $\OO_Y$, since its first
Chern class is zero. A nonzero section of the torsion
line $\omega_S$ trivializes it. Indeed its $m$th power
is a nonzero holomorphic function on the compact
connected reduced space and hence is a nonzero constant.
Thus, with $e=h^0(S,\omega_S)$, we have
\[
 e\in\{0,1\},\qquad
 e=1\Longleftrightarrow\omega_S\simeq\OO_S,
 \qquad h^1(S,\OO_S)=1+e.                         \tag{A.CD.2}
\]
Here $h^0(\OO_S)=1$ and Serre duality identifies
$h^2(\OO_S)$ with $e$.

Let $\mathcal Q=\eta_*\OO_T/\OO_S$ and let
$\Gamma\subset S$ be the conductor image scheme.
Both $S$ and $T$ are connected and compact, so their
global holomorphic functions are the same constants.
The normalization sequence consequently gives an
injection and the indicated kernel:
\[
 H^0(S,\mathcal Q)\hookrightarrow H^1(S,\OO_S),
 \qquad
 \operatorname{im}H^0(S,\mathcal Q)
 =\ker\bigl(H^1(S,\OO_S)\to H^1(T,\OO_T)\bigr).
                                                               \tag{A.CD.3}
\]
Because $H^2(Y,\Z)=0$, $\Pic(Y)$ is the connected
exponential image of $H^1(Y,\OO_Y)$. A nonconstant
holomorphic group homomorphism out of it has nonzero
differential. The differential of the stated Picard
pullback is the actual composite
\[
 H^1(Y,\OO_Y)\longrightarrow H^1(S,\OO_S)
 \longrightarrow H^1(T,\OO_T)
 \longrightarrow H^1(Z,\OO_Z).
\]
Thus its middle arrow has positive rank. Equations
(A.CD.2)--(A.CD.3) imply $h^0(\mathcal Q)\le e$.

The conductor is Cohen--Macaulay of dimension one,
and the exact duality comparison (A.RC.3) applies
without a reducedness assumption on $\Gamma$:
$\mathcal Q=\omega_\Gamma\otimes\omega_S^{-1}|_\Gamma$.
Serre duality on that curve therefore gives
\[
 h^1(\Gamma,\omega_S|_\Gamma)
       =h^0(S,\mathcal Q)\le e.                   \tag{A.CD.4}
\]
For any reduced irreducible component $D$ of $\Gamma$,
restriction of the line gives an exact quotient
$\omega_S|_\Gamma\to\omega_S|_D$. Its kernel is a
coherent sheaf on a one-dimensional compact space and
has no $H^2$, so the corresponding map on $H^1$ is
surjective. Hence
$h^1(D,\omega_S|_D)\le e$.

For $e=1$, (A.CD.2) makes this precisely
$p_a(D)=h^1(D,\OO_D)\le1$.
For $e=0$, the same inequality makes
$h^1(D,\omega_S|_D)=0$.
The restriction $\omega_S|_D$ is torsion by (A.CD.1)
and has degree zero. Curve Riemann--Roch now gives
\[
 0\le h^0(D,\omega_S|_D)
       =\chi(D,\omega_S|_D)=1-p_a(D),
\]
again proving $p_a(D)\le1$.
Finally, for a reduced irreducible curve the
normalization exact sequence gives
$p_a(D)=g(D^\nu)+\operatorname{length}
(\nu_*\OO_{D^\nu}/\OO_D)$.
Positive normalization genus therefore forces genus
one and zero normalization quotient. A normal complex
curve is smooth, which proves the final assertion.
\end{proof}

\begin{statement}{A.5.13}{Proved}{The conductor bound when the general fibre has $q=1$ and $p_g=0$}
Let $f:Y\to\PP^1$ be a proper connected-fibre map from
a smooth compact complex threefold with
$H^1(Y,\Z)=H^2(Y,\Z)=0$. Suppose its smooth general
fibre has $h^1(\OO)=1$ and $h^2(\OO)=0$.
Then
\[
 R^1f_*\OO_Y=\OO_{\PP^1},\qquad R^2f_*\OO_Y=0,
 \qquad\Pic(Y)=\C\beta,                            \tag{A.Q1.1}
\]
where the last identification uses the exponential
coordinates of D.1.9 and its nonzero transgression
class $\beta$. Every full fibre has structure-sheaf
cohomology dimensions $(1,1,0)$.
For an irreducible reduced support $S$, every
nonconstant ambient Picard pullback to a normalization
resolution forces every reduced conductor component
to have arithmetic genus at most one. In particular
each positive-genus component is smooth elliptic,
without any lci or conductor-reducedness hypothesis.
\end{statement}
\begin{proof}
By E4, $E=R^1f_*\OO_Y$ is a line bundle. The section
induced by D.1.9's integral transgression class is
nowhere zero: its image under the natural base-change
map at each point is the nonzero coefficient class on
that full fibre. A zero of the section in $E$ would
have zero image under that map. Thus $E\simeq\OO$.
Coherent Leray now gives $H^1(Y,\OO_Y)=\C\beta$.
The exponential sequence and the two stated integral
cohomology vanishings give the Picard identification
in (A.Q1.1).

The coherent sheaf $G=R^2f_*\OO_Y$ has generic rank
zero, hence is a torsion sheaf. Riemann--Roch on $Y$
gives $\chi(\OO_Y)=\int_Yc_1c_2/24=0$, since
$c_1(Y)$ is zero in rational cohomology. The Euler
characteristic of coherent Leray is therefore
\[
 0=\chi(\OO_{\PP^1})-\chi(E)+\chi(G)
      =\operatorname{length}G.
\]
Thus $G=0$. Top-degree base change gives
$H^2(F,\OO_F)=0$ for every full Cartier fibre $F$.
The multiplication-by-parameter sequence then gives
$H^1(F,\OO_F)=E\otimes\C(c)$, of dimension one;
E4 gives $h^0(F,\OO_F)=1$.

For a full fibre $F=mS$ with irreducible reduced
support, the quotient $\OO_F\to\OO_S$ induces a
surjection on $H^2$, because its kernel has support
of dimension two and has no $H^3$.
Hence $h^2(S,\OO_S)=0$. Riemann--Roch in the Cartier
sequence on $Y$ gives $\chi(\OO_S)=0$ as in A.5.12,
so $h^1(S,\OO_S)=1$.

Put $\mathcal Q=\eta_*\OO_T/\OO_S$ for the
normalization $T$ and write $\Gamma$ for the conductor
image scheme. Nonconstant Picard pullback has nonzero
differential factoring through
$H^1(S,\OO_S)\to H^1(T,\OO_T)$.
The source has dimension one, so this map has zero
kernel. The normalization sequence (A.CD.3) gives
$H^0(\mathcal Q)=0$. Conductor duality and Serre
duality give
\[
 H^1(\Gamma,\omega_S|_\Gamma)=0,
 \qquad H^1(D,\omega_S|_D)=0                      \tag{A.Q1.2}
\]
for each reduced irreducible component $D$, where
the second map is the restriction surjection in the
proof of A.5.12. Adjunction makes $\omega_S$ ambient,
so $\deg(\omega_S|_D)=0$ by $H^2(Y,\Z)=0$.
Curve Riemann--Roch yields
$0\le h^0(D,\omega_S|_D)=1-p_a(D)$.
The normalization genus and its nonnegative delta
length then give the asserted smooth elliptic
conclusion whenever that genus is positive.
\end{proof}

\begin{statement}{A.4.3}{Proved}{The omitted genus-one case for the standard rank-one fibre}
Let $E$ be an elliptic curve, let $P\in\Pic(E)$, and let $T=\PP_E(\OO_E\oplus P)$ with its zero and infinity sections $D_0,D_\infty$. Glue $D_0$ to $D_\infty$ by a translation of $E$ to form the compact normal-crossings surface $S$, and denote its normalization by $\eta:T\to S$. Every line bundle $A$ on $S$ whose pullback $\eta^*A$ is topologically trivial satisfies
\[
H^0(S,\widetilde\Omega_S^1\otimes A)=
\begin{cases}\C,&A\simeq\OO_S,\\0,&A\not\simeq\OO_S.\end{cases}
\tag{A.16}
\]
\end{statement}
\begin{proof}
For the ruling $p:T\to E$, every topologically trivial line bundle on $T$ is $p^*P_A$ for a unique $P_A\in\Pic^0(E)$. Indeed its degree on every ruling fibre is zero, so the adjunction map from the pullback of its rank-one direct image is an isomorphism, checked on each fibre by $H^0(\PP^1,\OO)=\C$ and $H^1(\PP^1,\OO)=0$. The line on $E$ has degree zero because its first Chern class pulls back to zero and $p$ has a section.

The relative cotangent sequence and restriction to ruling fibres give
\[
 H^0(T,\Omega_T^1\otimes p^*P_A)
 =p^*H^0(E,\Omega_E^1\otimes P_A).
\]
The relative quotient has no sections because its restriction to every fibre is $\OO_{\PP^1}(-2)$. Since $\Omega_E^1\simeq\OO_E$, the displayed group vanishes for nontrivial $P_A$. This proves (A.16) in that case by the injection into normalization differentials.

For $P_A=\OO_E$, choose a trivialization of $\eta^*A$. The descent datum for $A$ is a nonzero constant $\lambda$ identifying the two restrictions to $D_0,D_\infty$, because these are trivial lines on the compact connected curve $E$. Multiplying the chosen trivialization by a constant changes neither $\lambda$ nor its meaning. The local calculation for $xy=0$ identifies the image of $\widetilde\Omega_S^1$ in $\eta_*\Omega_T^1$ with pairs of forms whose tangential restrictions to the double curve agree. The unique normalization form is a scalar multiple of $p^*\alpha$, where $\alpha$ is a nonzero invariant form on $E$. Translation pulls $\alpha$ back to itself. Twisted descent is therefore the equation $\alpha=\lambda\alpha$, so a nonzero section exists exactly when $\lambda=1$. That equality also says precisely that the trivialization descends to $A\simeq\OO_S$. This proves both dimensions.
\end{proof}

\begin{statement}{A.4.4}{Proved}{Canonical triviality on an irreducible reduced torus fibre}
Let $g:Y\to C$ be a proper holomorphic map from a connected smooth compact threefold onto a smooth compact curve, with connected fibres and general fibre a complex two-torus. Every irreducible reduced full fibre $S=g^{-1}(c)$ has $\omega_S\simeq\OO_S$. When $H^2(Y,\Q)=0$, it therefore satisfies $h^q(S,\OO_S)=(1,2,1)$.
\end{statement}
\begin{proof}
The coherent sheaf $g_*\omega_Y$ is torsion-free: a nonzero function on an open subset of the base cannot annihilate a holomorphic section of a line bundle on the dominating smooth total space. It has generic rank one because the smooth torus fibre has one holomorphic top form; after restricting to a base open set where Grauert base change holds, this identifies its rank. On the smooth curve it is consequently a line bundle $M$. The evaluation morphism $g^*M\to\omega_Y$ is generically nonzero, so its zero scheme is an effective Cartier divisor $D$ and $\omega_Y\simeq g^*M\otimes\OO_Y(D)$. No component of $D$ dominates the base: on a general smooth torus fibre its defining section is a nonzero holomorphic top form and is nowhere zero. Thus $D$ is vertical.

All vertical components meeting the specified irreducible full fibre are that fibre itself. Consequently $\OO_Y(D)|_S$ is a power of $\OO_Y(S)|_S$, and the latter is trivial because a coordinate on $C$ defines the full reduced fibre. The pullback $g^*M|_S$ is also trivial. Adjunction gives $\omega_S=(\omega_Y\otimes\OO_Y(S))|_S\simeq\OO_S$. The fibre is reduced and connected, so $h^0(\OO_S)=1$; it is a Cartier divisor, hence Cohen--Macaulay and Gorenstein, so Serre duality gives $h^2(\OO_S)=1$. The divisor Riemann--Roch calculation in A.5.2 gives $\chi(\OO_S)=0$ under the stated rational cohomology hypothesis, yielding $h^1=2$.
\end{proof}

\begin{statement}{A.4.5}{Proved}{The full torsion-free cotangent assertion for irreducible normal-crossings torus fibres}
Let $g:Y\to C$ have the hypotheses of A.4.4 and let $S$ be a singular irreducible reduced full fibre with normal crossings. For every nontrivial unitary-flat line bundle $A$ on $S$,
\[
H^0(S,\widetilde\Omega_S^1\otimes A)=0,
\qquad
H^0(S,\Omega_Y^1|_S\otimes A)
 \simeq H^0(\widetilde S,\eta^*A).
\tag{A.17}
\]
The latter group has dimension one when $\eta^*A$ is holomorphically trivial, and zero otherwise. The normalization is exactly one of the following: a diagonal Hopf surface with two disjoint elliptic curves identified; a ruled surface over an elliptic curve with two disjoint sections identified by translation; or a rational surface with an anticanonical hexagon as in Schr\"oer--Siebert Proposition~2.6.
\end{statement}
\begin{proof}
We check the hypotheses of the classification, including its deformation sheaf. A.4.4 gives $\omega_S\simeq\OO_S$. Dualizing the locally free conormal resolution of $\Omega_S^1$ yields
\[
 \mathcal T^1_S=\mathcal Ext^1_{\OO_S}(\Omega_S^1,\OO_S)
 \simeq\OO_S/\mathcal J=\OO_{S_{\rm sing,red}}.
\]
Here the conormal is trivial because $S$ is the full reduced fibre; the local identities $\mathcal J=(x,y)$ at a double point and $\mathcal J=(xy,xz,yz)$ at a triple point show that this is the reduced singular locus. Hence $S$ is d-semistable in the exact sense of \cite[Definition~2.1]{SS2002}. That definition imposes irreducibility, normal crossings, d-semistability and trivial dualizing sheaf; it does not require a prior smoothing to Kodaira surfaces. Propositions~2.4, 2.5 and 2.6 there give the three normalization cases. In the ruled case their Proposition~4.1 says precisely that trivial $\omega_S$ forces the section identification to be a translation. The comparison with their notation is $X_{\mathrm{SS}}=S$ and $S_{\mathrm{SS}}=\widetilde S$.

In the rational case $\eta^*A$ is trivial: its first Chern class is torsion because the line is unitary-flat, whereas the integral second cohomology of a smooth rational surface is torsion-free and its Picard group injects into that group by the exponential sequence. Such a surface has no holomorphic one-forms, by birational invariance as proved in A.4.1. The injection $\widetilde\Omega_S^1\to\eta_*\Omega^1_{\widetilde S}$ therefore proves the first vanishing.

In the elliptic ruled case the proof of A.4.3 applies to a decomposable ruled surface of any degree: the two disjoint sections express it as $\PP_E(\OO_E\oplus P)$, with $P\in\Pic(E)$. The integral second cohomology of this ruled surface is torsion-free, by the projective-bundle formula. A unitary-flat line has torsion first Chern class, so its pullback is topologically trivial and comes from $\Pic^0(E)$. No degree assumption on $P$ occurs in the cotangent calculation, and A.4.3 proves the required vanishing for every nontrivial $A$.

In the Hopf case, pull back to a finite primary diagonal Hopf cover. Its universal covering is $\C^2\setminus\{0\}$ with contraction $(z_1,z_2)\mapsto(a_1z_1,a_2z_2)$, where $|a_i|<1$. A holomorphic one-form valued in a unitary-flat line lifts to $f_1(z)dz_1+f_2(z)dz_2$, with each $f_i$ holomorphic on the punctured space and hence on $\C^2$ by Hartogs' theorem. For a contraction character $\rho$ of modulus one, equivariance says that the coefficient of $z_1^{k_1}z_2^{k_2}dz_i$ can be nonzero only when
\[
 a_1^{k_1}a_2^{k_2}a_i=\rho
\]
(or its inverse, according to the convention for the associated line). The modulus on the left is strictly less than one and the character has modulus one. Thus every coefficient is zero. Finite etale pullback is injective on sections, so the twisted one-form group on the Hopf normalization also vanishes, and again the normalization injection proves the first assertion on $S$.

Finally A.1.2 and the first vanishing identify the ambient group with $H^0(\widetilde S,\eta^*A)$. A section of a unitary-flat line on a compact connected complex manifold is either zero or parallel and nowhere zero: its squared norm is locally a sum of squares of holomorphic functions, hence plurisubharmonic, and the maximum principle makes it constant. It therefore trivializes the line, and in the trivial case the section space is $\C$. This proves the dimension statement.
\end{proof}

\begin{statement}{A.4.6}{Proved}{Euler numbers of these irreducible normal-crossings fibres}
For $S$ in A.4.5,
\[
e(S)=\begin{cases}
2,&\widetilde S\text{ is rational},\\
0,&\widetilde S\text{ is elliptic ruled or Hopf}.
\end{cases}
\tag{A.18}
\]
Thus an arbitrary compact torus fibration whose fibres are irreducible reduced normal-crossings surfaces or multiples of smooth surfaces of Euler number zero has
\[
c_3(Y)=e(Y)=2r,
\tag{A.19}
\]
where $r$ is the number of rational-normalization fibres. For $b_1(Y)=b_2(Y)=0$, one has $r\in\{0,1\}$ and $b_3(Y)=2-2r$.
\end{statement}
\begin{proof}
Write $D\subset\widetilde S$ for the conductor support and $N\subset S$ for the singular curve. Normalization is a homeomorphism off these sets, so additivity of the topological Euler characteristic gives $e(S)=e(\widetilde S)-e(D)+e(N)$. In the Hopf case $e(\widetilde S)=0$, $D$ is a disjoint union of two elliptic curves, and $N$ is one elliptic curve; all three Euler numbers vanish. The ruled case has the same conductor curves and $e(\widetilde S)=e(\PP^1)e(E)=0$.

In the rational case the normalization is a Hirzebruch surface blown up twice, so its Euler number is $4+2=6$. The conductor is a cycle of six rational curves, with six nodes, so $e(D)=12-6=6$. Its image consists of three rational curves with exactly two triple points; counting the three preimages of each triple point once gives $e(N)=6-2(3-1)=2$. Hence $e(S)=2$.

Away from finitely many critical values the proper holomorphic map is a differentiable fibre bundle, and a two-torus has Euler number zero. Additivity for a stratification by this open base and the exceptional points therefore gives $e(Y)=\sum_c e((Y_c)_{\rm red})$. Nilpotents do not change the underlying topological space. The assumed fibre classes and (A.18) give (A.19). The equality with $c_3$ is the complex-oriented Gauss--Bonnet theorem. Poincare duality for the connected oriented six-manifold gives $e(Y)=2-2b_1+2b_2-b_3$. Under the two Betti vanishings this is $2-b_3$, so $b_3=2-2r\ge0$ and $r$ is zero or one.
\end{proof}

\begin{statement}{A.4.7}{Proved}{Unitary-flat differentials on every compact reduced complex surface}
Let $S$ be a compact reduced complex surface, let $\mathbb A$ be a rank-one unitary complex local system on its underlying topological space, and let $A=\mathbb A\otimes_{\C}\OO_S$ be its associated holomorphic line bundle. There is a canonical complex-linear injection
\[
 H^0(S,\widetilde\Omega_S^1\otimes A)
       \longrightarrow H^1(S,\mathbb A).
 \tag{A.20}
\]
The construction uses continuous local primitives that are holomorphic on the normalization; it does not assert that those primitives are holomorphic on a non-seminormal surface.
\end{statement}
\begin{proof}
Let $\eta:\widetilde S\to S$ be normalization and choose a resolution $\pi:Z\to\widetilde S$ such that the exceptional curves have normal crossings; the components of $Z$ are compact smooth surfaces. Pullback of Kähler differentials along $h=\eta\circ\pi:Z\to S$ kills the torsion of $\Omega_S^1$, because the resulting form vanishes on the dense inverse image of $S_{\rm reg}$ and the target is locally free. Thus a section $\alpha$ on the left of (A.20) gives a holomorphic one-form $\alpha_Z$ with values in the unitary-flat line $h^*A$. This map on sections is injective for the same dense-open reason.

Write $\nabla$ for the flat unitary connection. Holomorphicity gives $\nabla\alpha_Z$ type $(2,0)$. The Hermitian pairing identifies the tensor product of the flat line with its conjugate with the trivial local system. Stokes' theorem gives, on each connected component of $Z$,
\[
 \int_Z \nabla\alpha_Z\wedge\overline{\nabla\alpha_Z}
 =\int_Z d\bigl(\alpha_Z\wedge\overline{\nabla\alpha_Z}\bigr)=0.
\]
After multiplication by the fixed dimensional sign, the integrand is the nonnegative squared norm of a $(2,0)$-form times a positive volume form. Hence $\nabla\alpha_Z=0$.

We prove the required local primitive assertion at a point $x\in S$. Trivialize $\mathbb A$ on a sufficiently small connected neighbourhood $V$ of $x$; complex analytic spaces are locally contractible, so such a flat trivialization exists. The section $\alpha$ has a local lift to $\Omega_S^1\otimes A$, since the latter surjects onto the torsion-free quotient. Its pullback to each reduced fibre of $h$ over $x$ is zero: the morphism of that fibre to $S$ is constant, so the induced morphism on Kähler differentials is zero.

For each point $\widetilde x\in\eta^{-1}(x)$, a sufficiently small regular neighbourhood of the exceptional fibre $\pi^{-1}(\widetilde x)$ in the good resolution retracts onto that fibre. This neighbourhood can be constructed by plumbing the normal disc bundles of the exceptional curves at their transverse intersection points; it is a neighbourhood of the entire fibre. A sufficiently small inverse image of a neighbourhood of $\widetilde x$ lies inside it by properness. In the smooth case one instead uses a coordinate ball. Every period of $\alpha_Z$ on the regular neighbourhood is zero, because the form is closed and its restriction to the retracting fibre is zero. Integration from a fixed point therefore produces a single-valued holomorphic primitive: a smooth primitive has differential of type $(1,0)$, so its $\bar\partial$ is zero. Normality and $\pi_*\OO_Z=\OO_{\widetilde S}$ descend that primitive to a holomorphic function $b$ on a neighbourhood of $\widetilde x$. Choose its constant so that $b(\widetilde x)=0$. Perform this independently at every point over $x$ and shrink $V$ so that all its inverse image is covered by these neighbourhoods.

These primitives have equal values on every fibre of $\eta$ over $V$. To see this, form the finite analytic fibre product
$W=\widetilde S\times_S\widetilde S$ with projections $p_1,p_2$. On the normalization of every positive-dimensional irreducible component of $W_{\rm red}$, the differential of $p_1^*b-p_2^*b$ is zero on a dense open subset: both differentials are pullbacks of the same local lift of $\alpha$ from $S$. The isolated singularities of $\widetilde S$ can be omitted for this comparison; holomorphicity then extends the equality. Thus the difference is constant on each such component. The same conclusion is automatic on a zero-dimensional component. Shrink $V$ once more so that every component of the finite space $W|_V$ meets the fibre over $x$; components disjoint from that fibre can be removed by properness. Our normalization $b(\widetilde x)=0$ makes all these constants zero. Hence $p_1^*b=p_2^*b$ as continuous functions, proving equality on the fibres of $\eta$.

A finite proper surjection is a topological quotient map. The functions $b$ therefore descend to continuous functions on $V$, holomorphic after pullback to normalization, whose differential on the regular locus is $\alpha$. On overlaps, any two such primitives differ by a locally constant function: their difference is constant on each local normalized branch, and continuity identifies the constants whenever branches meet. With the flat transition functions of $\mathbb A$, these differences define a Cech class in $H^1(S,\mathbb A)$. Changing the chosen primitives changes it by a locally constant coboundary. This defines the canonical linear map (A.20).

If the resulting class is zero, adjust the local primitives by locally constant sections so that they glue. Their pullback is then a global holomorphic section of the unitary-flat line on each compact component of $Z$. Its squared norm is plurisubharmonic, so the maximum principle makes the section parallel; its covariant differential is zero. Consequently $\alpha_Z=0$, and injectivity of pullback gives $\alpha=0$. This proves (A.20).
\end{proof}

\begin{statement}{A.4.8}{Proved}{The torsion-free cotangent assertion for arbitrary singular supports in a torus fibration}
Let $g:Y\to\PP^1$ be a proper surjective holomorphic map with connected fibres from a smooth compact complex threefold with $b_1(Y)=b_2(Y)=0$, and suppose its general fibre is a complex two-torus. Let $L_r=\exp(2\pi\mathrm i r\beta)\in\Pic^0(Y)$ be the transgression family constructed in D.1.9. For every real irrational $r$ and every reduced fibre support $S=(Y_c)_{\rm red}$,
\[
 H^1(S,\mathbb A_{r,S})=0,
 \qquad
 H^0(S,\widetilde\Omega_S^1\otimes L_r|_S)=0,
 \tag{A.21}
\]
where $\mathbb A_{r,S}$ is the unitary local system whose associated holomorphic line is $L_r|_S$. In particular this proves Equation~(11) of CDP20 for all singular reduced fibre supports in the torus case, with no assumption on their singularities or their scheme-theoretic multiplicities.
\end{statement}
\begin{proof}
We specify all comparison maps used from D.1.9. On a sufficiently small disc $\Delta$ about $c$, put $U=g^{-1}(\Delta)$ and $U^*=U\setminus S$. The integral section $\gamma$ of $R^1g_*\Z$ has a lift $\gamma_U\in H^1(U,\Z)$ by Leray over the contractible disc. Its image in $H^1(U,\OO_U)$ equals $\beta|_U$, by the coherent Leray sequence. Thus the restriction of the holomorphic line $L_r$ to $U$ is the line associated to the flat character
\[
 \pi_1(U)\longrightarrow S^1,\qquad
 a\longmapsto\exp\bigl(2\pi\mathrm i r\langle\gamma_U,a\rangle\bigr).
 \tag{A.22}
\]
Denote this local system by $\mathbb A_U$. D.1.8--D.1.9 prove that $\gamma_U$ restricts nontrivially to every nearby smooth fibre. Irrationality of $r$ therefore makes (A.22) nontrivial on such a fibre $F$.

A nontrivial rank-one local system on the real torus underlying $F$ has zero cohomology in every degree. Explicitly, use the standard cellular cochain complex of the product of four circles with coefficients in the character. It is the Koszul complex for the four scalars $\chi(e_j)-1$. At least one scalar is nonzero; contraction in that factor, divided by this scalar, is a contracting homotopy. Therefore $H^q(F,\mathbb A_U|_F)=0$ for every $q$. The map $U^*\to\Delta^*$ is a differentiable fibre bundle. Its local-system Leray spectral sequence has zero fibre-cohomology sheaves, hence
$H^1(U^*,\mathbb A_U|_{U^*})=0$.

Inclusion induces a surjection $\pi_1(U^*)\to\pi_1(U)$: a loop in the smooth real six-manifold $U$ can be perturbed off the divisor, whose real codimension is at least two, relative to a base point in $U^*$. The induced map on first cohomology with these local coefficients is injective. For an explicit verification, identify first cohomology with crossed homomorphisms modulo coboundaries. If a crossed homomorphism $a:G\to M$ becomes $a(q(h))=(\chi(q(h))-1)v$ after pullback by a surjection $q:H\to G$, surjectivity gives $a(g)=(\chi(g)-1)v$ for every $g\in G$. Thus its original class was already zero. Applied to the displayed fundamental-group map, this proves $H^1(U,\mathbb A_U)=0$.

The argument applies after every sufficiently small shrinking of $\Delta$, with the restrictions of the same local system. Taking the direct limit and using proper base change identifies the resulting stalk with $H^1(S,\mathbb A_U|_S)$; scheme-theoretic nilpotents do not affect this topological local system. This proves the first vanishing in (A.21). The second is exactly A.4.7 applied to the restriction of (A.22). All singular values are treated by the same global parameter $r$, so the assertion is simultaneous.
\end{proof}

\begin{statement}{A.4.9}{Proved}{CDP20 Proposition 7.3 with a complete replacement proof}
Let $Y$ be a connected smooth compact complex threefold with $H^1(Y,\Z)=H^2(Y,\Z)=0$, algebraic dimension one and holomorphic algebraic reduction $g:Y\to C$. There is a dense open subset of $\Pic(Y)$, whose complement is a proper complex analytic subset, such that
\[
 H^0\bigl((Y_c)_{\rm red},
       \widetilde\Omega^1_{(Y_c)_{\rm red}}\otimes L|_{(Y_c)_{\rm red}}\bigr)=0
 \tag{A.23}
\]
for every singular reduced fibre support and every $L$ in that subset. In fact the transgression twists $L_r$ of D.1.9 satisfy (A.23) on every reduced fibre support, including smooth and reducible supports, for every real irrational $r$.
\end{statement}
\begin{proof}
The independently rederived smooth-fibre classification C.3.1 leaves complex tori, bielliptic surfaces, primary or secondary Kodaira surfaces, and primary or secondary Hopf surfaces or Inoue surfaces. We prove that for each such smooth fibre $F$ and each nonzero $\gamma_F\in H^1(F,\Z)$ the local system
\[
 \mathbb A_{r,F}:\quad a\longmapsto
       \exp(2\pi\mathrm i r\langle\gamma_F,a\rangle)
 \quad(r\in\mathbb R\setminus\Q)
\]
satisfies $H^0(F,\mathbb A_{r,F})=H^1(F,\mathbb A_{r,F})=0$. Its character has infinite image, so the assertion in degree zero is immediate.

For a torus the Koszul argument in A.4.8 applies. A bielliptic surface has a finite étale torus cover. Pullback of the integral class is nonzero: transfer composed with pullback is multiplication by the degree. The character stays nontrivial after pullback, and transfer with local coefficients makes pullback on cohomology injective over $\C$. Thus degree-one vanishing descends from the torus cover.

A primary Kodaira surface is a principal elliptic bundle over an elliptic curve, as described in \cite[\S1 and the central-extension diagram in the proof of Lemma~1.1]{SS2002}. Its fundamental group has a central extension
\[
 1\longrightarrow\Z^2\longrightarrow\Gamma\longrightarrow\Z^2
       \longrightarrow1.
 \tag{A.24}
\]
For any nontrivial one-dimensional character $\chi$ of $\Gamma$, if its restriction to the central subgroup is nontrivial, the latter's cohomology with these coefficients is zero by the two-variable Koszul contraction. The Hochschild--Serre spectral sequence then gives $H^1(\Gamma,\C_\chi)=0$. If its central restriction is trivial, $\chi$ descends to a nontrivial character of the quotient $\Z^2$. Centrality makes its action on the cohomology of the central subgroup equal to that scalar character, and the same Koszul contraction kills every term on the second page. Again $H^1(\Gamma,\C_\chi)=0$. First cohomology of a connected space with local coefficients is first cohomology of its fundamental group, so this proves the required assertion. A secondary Kodaira surface has a finite primary cover; the same transfer argument as above applies. Hopf surfaces have finite primary covers, and a primary Hopf surface is diffeomorphic to $S^1\times S^3$; the cover statement and topology are recorded in \cite[\S2.1.1 and Theorem~2.2]{Albanese2020}. The nontrivial circle character has an invertible cellular differential, giving the required vanishing and its finite-cover descent.

For completeness we also give the Inoue calculation, rather than assuming a generic-vanishing theorem for its cotangent bundle. The explicit constructions in \cite[\S2.1.2 and Proposition~2.3]{Albanese2020} identify every Inoue surface with a mapping torus $N_\varphi$ of a compact three-manifold $N$. In the $S_M$ cases $N$ is a three-torus and the eigenvalues of $\varphi^*$ on $H^1(N,\C)$ are $\alpha,\beta,\bar\beta$, with $\alpha>1$ and $|\beta|=\alpha^{-1/2}<1$, after reversing the circle orientation when necessary. In the $S^+$ and $S^-$ cases $N$ is the nontrivial circle bundle over a two-torus; the nonzero Euler class makes $H^1(N,\C)$ the pullback of $H^1(T^2,\C)$ by its Gysin sequence. The monodromy on it has eigenvalues $\alpha,\alpha^{-1}$ or $\alpha,-\alpha^{-1}$, respectively. These assertions follow as well directly from the conjugation relations displayed in that construction: on the abelian rank-three kernel the matrix is $M$, and on the quotient of the Heisenberg kernel by its centre the matrix is $N\in\mathrm{GL}_2(\Z)$ with precisely those eigenvalues.

None of these degree-one eigenvalues has modulus one. In particular there are no invariant degree-one classes on $N$. The ordinary mapping-torus sequence therefore identifies $H^1(F,\Z)$ with the copy of $\Z$ from the circle, since the invariant subgroup of the torsion-free group $H^1(N,\Z)$ is zero. Thus $\mathbb A_{r,F}$ is pulled back from the circle, with monodromy $\lambda\in S^1$ of infinite order. The local-coefficient Leray sequence for $N\to N_\varphi\to S^1$ gives
\[
\begin{split}
0\longrightarrow
 \coker(\lambda-1:H^0(N,\C)\to H^0(N,\C))
 \longrightarrow H^1(F,\mathbb A_{r,F})\\
 \longrightarrow
 \ker(\lambda\varphi^*-1:H^1(N,\C)\to H^1(N,\C))
 \longrightarrow0.
\end{split}
\]
The first term is zero because $\lambda\ne1$, and the last term is zero by the eigenvalue moduli. Using the reciprocal convention for the circle local system replaces $\lambda$ by $\lambda^{-1}$ and gives the same conclusion. This completes the degree-one calculation for every surface type in C.3.1.

Now repeat the neighbourhood argument of A.4.8 with these two fibre-cohomology vanishings. D.1.8--D.1.9 supply the same nonzero integral class on every nearby smooth fibre and the unitary local system on every fibre neighbourhood. The only terms that can contribute to first cohomology of the punctured neighbourhood are its degree-zero and degree-one fibre cohomology, just proved zero. The fundamental-group surjection and proper base change therefore give $H^1((Y_c)_{\rm red},\mathbb A_{r,c})=0$ for every $c$. Applying A.4.7 proves (A.23) for the single global twist $L_r$, simultaneously on all reduced supports.

Finally, the exponential sequence and the integral vanishings identify $\Pic(Y)$ with the complex vector space $H^1(Y,\OO_Y)$. Choose holomorphic Cech cocycles representing a basis; exponentiating their linear combinations gives a holomorphic family of line bundles on $Y\times H^1(Y,\OO_Y)$ whose fibre is the indicated line bundle. There are finitely many singular fibre supports. For each of them, tensor its fixed coherent torsion-free cotangent sheaf by this family and apply proper semicontinuity: the locus where its space of sections is nonzero is a closed complex analytic subset of the parameter space. Each is proper because it excludes $L_r$. Their finite union is proper and has empty interior, giving the asserted dense open set. This proves the statement of Proposition~7.3. The invalid genus-two deduction in its printed proof remains refuted by A.4.2; the proof above replaces that argument through the explicit maps (A.20)--(A.24).
\end{proof}
